\documentclass[
  aps,            % American Physical Society style
  prc,            % Physical Review C
  reprint,        % two-column published format; use "preprint" for the one-column manuscript style
  superscriptaddress,
  floatfix,
  nofootinbib
]{revtex4-2}

% ---------------------------------------------------------------------------
%  Hybrid resolution refinement + rodeo algorithm on the pairing Hamiltonian
%  Built from the Jupyter notebook ResolutionRefPairing.ipynb
%  and Qian et al., Eur. Phys. J. A 60, 151 (2024).
%  RevTeX 4-2 / APS style.
% ---------------------------------------------------------------------------

\usepackage{amsmath,amssymb}
\usepackage{physics}
\usepackage{graphicx}
\usepackage{quantikz}   % circuit diagrams (Figs. \ref{fig:circ-vqe}, \ref{fig:circ-rodeo})
\usepackage{hyperref}

\newcommand{\Hlow}{H_{\mathrm{low}}}
\newcommand{\Hhigh}{H_{\mathrm{high}}}
\newcommand{\Eobj}{E_{\mathrm{obj}}}

\begin{document}

\title{Resolution refinement followed by the rodeo algorithm for
eigenstate preparation: implementation and assessment on the pairing Hamiltonian}

\author{Scott Bogner}
\affiliation{Facility for Rare Isotope Beams and Department of Physics and Astronomy, Michigan State University, East Lansing, MI 48824, USA}
\author{Cecilie Glittum}
\affiliation{Department of Physics, University of Oslo, N-0316 Oslo, Norway}
\author{Heiko Hergert}
\affiliation{Facility for Rare Isotope Beams and Department of Physics and Astronomy, Michigan State University, East Lansing, MI 48824, USA}
\author{Morten Hjorth-Jensen}
\email{mhjensen@uio.no}
\affiliation{Department of Physics, University of Oslo, N-0316 Oslo, Norway}
\author{Gunnar Lange}
\affiliation{Department of Physics, University of Oslo, N-0316 Oslo, Norway}
\author{Ryan LaRose}
\affiliation{Department of Computational Mathematics, Science and Engineering and Facility for Rare Isotope Beams, Michigan State University, East Lansing, MI 48824, USA}
\author{Dean Lee}
\affiliation{Facility for Rare Isotope Beams and Department of Physics and Astronomy, Michigan State University, East Lansing, MI 48824, USA}
\author{Francesco Pietro Massel}
\affiliation{Department of Science and Industry Systems, University of South-Eastern Norway, N-3616 Kongsberg, Norway}
\affiliation{Department of Physics, University of Oslo, N-0316 Oslo, Norway}
\author{Jean Paul Sadia}
\affiliation{Department of Computational Mathematics, Science, and Engineering, Michigan State University, East Lansing, MI 48823, USA}
\affiliation{Center for Quantum Computing, Science, and Engineering, Michigan State University, East Lansing, MI 48823, USA}
\date{\today}

\begin{abstract}
Projective spectral filters such as the rodeo algorithm prepare eigenstates with
exponential accuracy in the number of filter cycles, but their cost is set by
the squared overlap $p$ of the input state with the target.  We implement and
assess, on the exactly solvable constant-pairing Hamiltonian, a pipeline in
which \emph{resolution refinement} --- the adiabatic transport of an eigenstate
from a coarse to a fine single-particle basis along a shift-controlled,
gapped path --- supplies the input of the rodeo filter, and we compare it
against the simplest alternatives on the same footing: the same filter applied
to the unrefined prolonged state and to a full-space unitary coupled-cluster
(UCCSD) state.  All comparisons use the ground-state infidelity as the common
observable, the same second-order product-formula evolution, and a two-qubit
gate budget per accepted preparation that includes the embedded
coarse-Hamiltonian term of the refinement path.  For four particles on four
levels (eight system qubits) the refined input ($p=0.998$) is filtered to an
infidelity of $5\times10^{-4}$ after two cycles and $1.5\times10^{-4}$ after
six, close to the floor set by the Trotterized evolution; the same filter brings
the prolonged determinant ($p=0.868$) to a comparable accuracy at a comparable
gate count, while the UCCSD state, whose actual infidelity is $1.8\times10^{-4}$
although its energy error is $1.4\times10^{-3}$, reaches $1.3\times10^{-4}$
after two cycles at one third of the cost.  Refinement therefore does not
beat a cheap, accurate variational input at this system size; its merits are
the absence of a classical optimization loop and of training measurements.
The extracted refinement time grows with the inverse of the minimum path gap
when the gap is varied through the energy shift, and with the loss of initial
overlap when the pairing strength is varied.  A pair-breaking particle-hole
interaction, which destroys the seniority symmetry of the model, leaves the
structure and convergence of the pipeline unchanged, raises the depth of a
Trotter step ninefold, and leaves the hybrid tied with the filtered UCCSD state
at weak and behind it at strong pair breaking.  Level occupation numbers and
correlation energies of the prepared states are reported, and all stages are
compiled to rotation--CNOT circuits with gate counts verified against exact
evolution.
\end{abstract}

\maketitle



% ===========================================================================
\section{Introduction}
% ===========================================================================
Determining the low-lying eigenvalues and eigenvectors of a many-body Hamiltonian is
one of the grand challenges of quantum many-body theory, where the Hilbert-space
dimension grows exponentially with the number of constituents.  Quantum computers
offer new routes to eigenvalue estimation and eigenstate preparation, and their
prospects for nuclear structure have been mapped in recent community
studies~\cite{Beck2023,Bauer2023,LeeReview2023}, with schematic nuclear models ---
pairing, Lipkin, and shell-model Hamiltonians --- serving as the standard proving
ground~\cite{Ovrum2007,Hlatshwayo2022,Stetcu2022,PerezObiol2023,Jiang2023,Zhang2024,Lacroix2023}.
The state-of-the-art preparation methods, however, each face an obstacle at scale.
The variational quantum eigensolver (VQE)~\cite{Peruzzo2014} is limited by the cost
of its classical optimization loop, by the accuracy floor of its ansatz and, for
sufficiently expressive circuits, by barren plateaus --- gradients that vanish
exponentially with system size~\cite{McClean2018,Larocca2025}.  Adiabatic state
preparation requires evolution times that diverge if the gap between the tracked
state and the rest of the spectrum closes along the path~\cite{Albash2018}.
Projective spectral filters --- quantum phase estimation~\cite{Abrams1999},
polynomial filters based on quantum signal processing~\cite{LinTong2020,Dong2022},
projection by measurement~\cite{Stetcu2023}, and the rodeo algorithm
(RA)~\cite{Choi2021,Qian2024,Cohen2023,BeeLindgren2022,Patkowski2026} --- converge
exponentially in the number of cycles, but only once the initial state has a
significant squared overlap $p$ with the target eigenstate; their cost grows as
$1/p$, and for a fixed, unoptimized reference state $p$ falls exponentially with
system size.

These limitations are complementary, and the idea of \emph{preconditioning} the
input of a projective filter by an adiabatic stage is as old as the rodeo
algorithm itself: the original proposal~\cite{Choi2021} already combines a
short adiabatic ramp with the filter, the hierarchical fusion method of
Ref.~\cite{Patkowski2025fusion} builds large systems from exactly solvable
subsystems by alternating adiabatic ramps and rodeo cycles, and the
resolution-refinement proposal~\cite{Bogner2026} explicitly recommends
switching from adiabatic evolution to a filter such as the rodeo algorithm once
the overlap is large.  \emph{Resolution refinement} bootstraps an eigenstate
across model-space scales: an eigenstate of a coarse (low-resolution)
Hamiltonian is lifted by a prolongation operator into a finer space and
adiabatically evolved to the fine-resolution eigenstate.  Because refinement
changes the low-energy spectrum only mildly, the interpolation path can be
kept gapped, and under the additional assumption that the first-order
perturbative correction to the tracked state stays bounded~\cite{Bogner2026}
the required adiabatic time scales as the inverse gap times a slowly growing
system-size factor rather than as the inverse gap squared.  The present work
implements this chain --- (i) prepare a low-resolution eigenstate; (ii) prolong
and adiabatically refine for a time $T$; (iii) apply the rodeo algorithm ---
for a hierarchy of fermionic single-particle bases, and asks the question that
decides whether the refinement stage earns its gates: does refinement followed
by filtering beat filtering the best input that is available cheaply?

The pairing problem has become a standard target for quantum algorithms, and
the programme of Lacroix and collaborators is the closest in spirit to ours.
Its common thread is the deliberate use of symmetry breaking and
restoration: a symmetry-breaking (BCS-like) trial state is cheap and
correlated, and the symmetry is restored afterwards by a projective step ---
quantum phase estimation used as a projector~\cite{Lacroix2020}, Grover-type
oracles~\cite{RuizGuzman2023}, spin filters~\cite{Siwach2021} or classical
shadows~\cite{RuizGuzman2024} --- with the restored state then handed to a
spectral post-processing stage such as phase estimation or a Krylov
expansion~\cite{RuizGuzman2022,RuizGuzman2021}; Refs.~\cite{Lacroix2023EPJA,Lacroix2023}
review this viewpoint and Ref.~\cite{Ayral2023} the broader many-body landscape.  More
recently the same group has used ADAPT-VQE on pairing Hamiltonians with
symmetry-conserving and symmetry-breaking operator pools~\cite{Zhang2024},
extracted excited states from the subspace spanned along the ADAPT convergence
path~\cite{Zhang2025}, obtained Lipkin spectra with a quantum
equation-of-motion approach~\cite{Hlatshwayo2022} and with a two-qubit
generator-coordinate scheme~\cite{BeaujeaultTaudiere2024}, and
built one-body Green's functions of small superfluid systems from a
variational even-$N$ ground state and a quantum subspace expansion for its
$N\pm1$ neighbors~\cite{AychetClaisse2026}.  Our pipeline addresses the step
that all of these share --- obtaining a high-fidelity eigenstate of the
fine-resolution Hamiltonian from a cheap approximate one --- with a different
mechanism: not symmetry projection or subspace diagonalization, but an
adiabatic path across resolutions followed by an energy filter; the relation
is taken up in Sec.~\ref{sec:discussion}.

In this work we realize the full pipeline end to end on a transparent and exactly
solvable testbed, the constant-pairing Hamiltonian familiar from nuclear and
condensed-matter many-body theory~\cite{LNP936} and integrable by the Richardson
ansatz~\cite{Richardson1963,RichardsonSherman1964,Dukelsky2004}.  Its exact
solvability is precisely what makes it the right laboratory: every stage of the
pipeline can be validated quantitatively --- the gate-level UCCSD coarse solver
against full configuration interaction (FCI) and coupled-cluster doubles (CCD); the
prolongation and adiabatic refinement against the exact high-resolution eigenstate;
and the rodeo filter against the analytic success-probability formulas of
Ref.~\cite{Qian2024}.  All circuits are implemented explicitly as rotation and CNOT
gates (and an ancilla-based rodeo cycle), so the demonstration doubles as a
reference implementation for hardware.

Relative to the resolution-refinement proposal of Ref.~\cite{Bogner2026}, the
rodeo literature~\cite{Choi2021,Qian2024} and the fusion
method~\cite{Patkowski2025fusion}, which established adiabatic preconditioning
of the filter on lattice and spin systems, the present work contributes five
elements: (i) an implementation of refinement for a hierarchy of
\emph{single-particle bases} in a fermionic pairing problem, where the
prolongation is gate-free and the embedded coarse Hamiltonian is a conditioned
operator whose circuit we give and count; (ii) a quantitative benchmark of the
refinement-plus-filter chain against the same filter applied to the unrefined
input and to a full-space UCCSD state, with the ground-state infidelity as the
common observable and matched evolution accuracy; (iii) a gate-level
compilation of every stage --- Trotterized refinement and Trotterized
controlled evolution included --- with a two-qubit-gate budget per accepted
preparation that separates one-time training, repeated preparation and the
filter; (iv) a test of the pipeline on a seniority-breaking Hamiltonian --- the
pairing model supplemented by a particle-hole interaction~\cite{MHJbook}, which
carries the one-particle--one-hole channel of realistic interactions ---
including the corresponding product-formula study and resource table; and (v)
a reproducible reference implementation in which every figure and number
derives from a single executable notebook.  The outcome of (ii) is stated
plainly: at the system sizes accessible to exact validation here, refinement
followed by filtering does not outperform filtering a UCCSD input, and is
roughly tied with filtering the unrefined prolonged state at matched fidelity;
the case for the refinement stage rests on its freedom from classical
optimization and training measurements, and on regimes --- larger spaces,
poorer variational inputs --- that lie beyond what we can validate exactly.

% ===========================================================================
\section{The pairing model and its qubit encoding}
% ===========================================================================
We study the constant-pairing Hamiltonian
\begin{widetext}
\begin{equation}
  H=\delta\sum_{p=0}^{k-1} p\sum_{\sigma=\uparrow\downarrow} n_{p\sigma}
   -\frac{g}{2}\sum_{p,q=0}^{k-1} a^\dagger_{p\uparrow}a^\dagger_{p\downarrow}
     a_{q\downarrow}a_{q\uparrow},
  \label{eq:H}
\end{equation}
\end{widetext}
with $k$ doubly degenerate single-particle levels of spacing $\delta$ (set to
$\delta=1$) and a pairing strength $g$.  The interaction scatters time-reversed
pairs between levels; it conserves the seniority (the number of unpaired particles),
so the ground state of $N$ paired particles lives entirely in the seniority-zero
subspace, and the model is exactly solvable by the Richardson
ansatz~\cite{Richardson1963,RichardsonSherman1964,Dukelsky2004} --- the property
that anchors every benchmark below.  We map level $p$ and spin $\sigma$ to qubit $2p+\sigma$ (so a $k$-level
problem uses $2k$ qubits) and encode the fermionic operators through the
Jordan--Wigner transformation.  The Hartree--Fock (HF) reference fills the lowest
$N/2$ levels, with $E_{\rm HF}(N{=}4,g)=2-g$.

Throughout we use $N=4$ particles as the running example.  The fixed-$N$ sector
dimension is $\binom{2k}{N}$, growing from $1$ at $k=2$ (a single filled
determinant) to $70$ at $k=4$ and $1820$ at $k=8$; the seniority-zero ground-state
problem is correspondingly compact, which makes the model an ideal multi-resolution
laboratory.  ``Basis refinement'' here means \emph{increasing $k$ at fixed $N$}: the
number of single-particle levels (and qubits) grows while the particle number is
held fixed.

% ===========================================================================
\section{The hybrid pipeline}
% ===========================================================================

\subsection{Coarse-state preparation: gate-level UCCSD--VQE}
\label{sec:uccsd}
The coarse eigenstate is prepared with a unitary coupled-cluster ansatz restricted
to singles and doubles (UCCSD)~\cite{LeeUCC2019,Tilly2022}, Trotterized into a
product of excitation circuits acting on the HF reference,
\begin{equation}
  \ket{\Psi(\boldsymbol\theta)}
  =\prod_{\mu} e^{\theta_\mu(\hat\tau_\mu-\hat\tau_\mu^\dagger)}\ket{\rm HF}.
\end{equation}
Through the Jordan--Wigner map each anti-Hermitian generator becomes a sum of
commuting Pauli strings, so every factor compiles into a basis change, a CNOT ladder,
a single $R_z(\theta)$ rotation, and the uncompute; a single excitation yields two
Pauli-rotation sub-circuits and a double excitation yields eight (the circuits are
shown in Fig.~\ref{fig:circ-vqe}).  We optimize with
the exact analytic gradient of the product ansatz (the naive two-term parameter shift
is incorrect beyond the first factor, because later factors acquire frequency-one
cross terms).  Because the pairing force does not couple to seniority-changing
single excitations, all singles amplitudes vanish to machine precision and the
solver is effectively a \emph{doubles} theory, mirroring the classical CCD structure.

\subsection{The prolongation operator}
The prolongation $P$ embeds a low-resolution state into the high-resolution space.
For basis refinement it maps every many-body basis state of the $k_{\rm low}$ space to
the \emph{same} occupation pattern in the $k_{\rm high}$ space, leaving the new
single-particle levels empty.  In canonical Jordan--Wigner ordering this is
gate-free; for lattice refinement the analogous operator distributes each occupied
coarse site over the neighboring fine sites with two single-qubit rotations and two
CNOTs per site per spatial direction, i.e.\ $O(N_{\rm sites})$ gates per
dimension~\cite{Bogner2026}.  The prolonged state $\ket{\Phi_0}=P\ket{\psi_{\rm low}}$
satisfies $\expval{\Hhigh}{\Phi_0}=E_{\rm low}$ exactly, and already provides a
high-quality first guess for the high-resolution ground state.

\subsection{Resolution refinement (adiabatic phase)}
\label{sec:refine}
The prolonged state is refined by adiabatically following the interpolated,
energy-shifted Hamiltonian
\begin{widetext}
\begin{equation}
  H(s)=\cos^2\!\Big(\tfrac{\pi s}{2}\Big)\,P(\Hlow-\mu)P^\dagger
      +\sin^2\!\Big(\tfrac{\pi s}{2}\Big)\,(\Hhigh-\mu),
  \qquad s=\frac{t}{T},
  \label{eq:path}
\end{equation}
\end{widetext}
evolved by a Suzuki--Trotter product over total time $T$.  The shift $\mu$
(here $\mu=E_{\rm low}+0.6$) places the tracked low-energy states below the large
null space of $P^\dagger$, so the level followed from $s=0$ to $s=1$ is the
high-resolution ground state.  The first term is \emph{not} the coarse
Hamiltonian acting on the coarse register alone: since $PP^\dagger$ projects on
the states in which the appended single-particle levels are empty,
\begin{equation}
  A\equiv P(\Hlow-\mu)P^\dagger=(\Hlow-\mu)\otimes
  \ket{0\!\cdots\!0}\!\bra{0\!\cdots\!0}_{\rm app},
  \label{eq:Aterm}
\end{equation}
a coarse-space operator conditioned on the $2(k_{\rm high}-k_{\rm low})$ new
qubits being in $\ket{0}$; its circuit and cost are given in
Sec.~\ref{sec:trotterized}.  At $s=0$ the gap above the tracked state is
$\min(E_1^{\rm low}-E_0^{\rm low},\,\mu-E_0^{\rm low})$: increasing the shift
lowers the embedded coarse states further below the null space, and the gain
in initial gap saturates at the coarse excitation gap (for a one-dimensional
coarse space, as in the $k=2$ start used below, it is simply $\mu-E_{\rm low}$).
Along the path the instantaneous spectrum is found to stay smooth and gapped
in all cases studied here.  A gapped path alone does not fix the adiabatic
time scale; Ref.~\cite{Bogner2026} shows that when the squared norm of the
first-order perturbative correction to the tracked state stays bounded along
the path --- a condition on the transition matrix elements of $dH/ds$ relative
to the excitation gaps --- the adiabatic time scales as $T\sim1/\Delta E$
rather than as the $O(1/\Delta E_{\min}^2)$ of the generic adiabatic
theorem~\cite{Albash2018}.  We test this empirically in
Sec.~\ref{sec:res-gscan} and treat the overlap
$p=\abs{\braket{\Psi_{\rm high}}{\Phi(T)}}^2$ as the figure of merit of the
stage.

\subsection{The rodeo algorithm}
\label{sec:rodeo}
The rodeo algorithm~\cite{Choi2021,Qian2024} removes the residual error left by the
adiabatic tail.  One cycle acts on the system register plus a single ancilla
(reused across cycles by mid-circuit measurement; the circuit is shown in
Fig.~\ref{fig:circ-rodeo}): a Hadamard on the ancilla, the
controlled evolution $e^{-iH t_k}$ of the system, a phase gate $P(E\,t_k)$ on the
ancilla (which effectively shifts $H\to H-E$), a second Hadamard, and an ancilla
measurement.  Post-selecting the ancilla in $\ket{0}$ applies the non-unitary filter
$\tfrac12\big(\mathbb{1}+e^{iE t_k}e^{-iH t_k}\big)$ to the system: an eigenstate
$\ket{E_j}$ acquires the factor $e^{i(E-E_j)t_k/2}\cos[(E_j-E)t_k/2]$.  For an
eigenstate input the probability that all $M$ ancilla measurements return $0$ is
therefore~\cite{Qian2024}
\begin{equation}
  P_{0^{M}}\!\big(E\mid\{t_k\}\big)=\prod_{k=1}^{M}
  \cos^2\!\Big[\tfrac{t_k}{2}\big(\Eobj-E\big)\Big],
  \label{eq:rodeo1}
\end{equation}
and averaging the $t_k$ over a Gaussian of RMS width $\sigma$ gives
\begin{equation}
  P_{0^{M}}(E)=\Bigg[\frac{1+e^{-(\Eobj-E)^2\sigma^2/2}}{2}\Bigg]^{M}.
  \label{eq:rodeo2}
\end{equation}
On resonance ($E=\Eobj$) the success probability is unity at every cycle, while
off-resonant weight is suppressed geometrically toward the background $2^{-M}$ (the
arithmetic mean of $\cos^2$ is $1/2$); note that the Gaussian-averaged
single-cycle factor $q=\tfrac12[1+e^{-(\Eobj-E)^2\sigma^2/2}]$ is never
smaller than $1/2$, so $2^{-M}$ is the suppression of components far from
resonance, $\abs{\Eobj-E}\sigma\gg1$, not a general lower bound on the
retained weight.  Scanning $E$ thus locates eigenvalues as
peaks whose width $\sim 1/\sigma$ sets the energy resolution; fixing $E$ at the
target eigenvalue turns the algorithm into a state-preparation filter.

The controlled evolution is the only system-size-dependent ingredient.  On hardware
$e^{-iH t}$ is realized by a Suzuki--Trotter decomposition over the $2k$
qubits~\cite{Trotter1959,Suzuki1976,Childs2021}; for the single-qubit benchmark of
Ref.~\cite{Qian2024} it reduces to a $Z$--$Y$--$Z$ Euler rotation
$R_{\hat n}(\theta)=e^{-i\theta\,\hat n\cdot\vec\sigma/2}$ implemented through the
Qiskit native gate.  In Secs.~\ref{sec:res-uccsd}--\ref{sec:res-rodeo} we apply the
controlled evolution exactly within the conserved $N$-particle sector to isolate the
algorithmic behavior; the circuit-level pipeline with explicit product formulas,
and the Trotter error it introduces, are quantified in Sec.~\ref{sec:trotterized}.
Throughout we draw the cycle times $t_k$ from a Gaussian as in
Ref.~\cite{Choi2021}; optimized deterministic~\cite{Cohen2023} or
geometric~\cite{Patkowski2026} time schedules suppress off-resonant weight faster
at fixed runtime and would only improve the numbers reported here, as would the
controlled-reversal-gate compilation of the controlled evolution of
Ref.~\cite{BeeLindgren2022}.

\paragraph{Convergence criterion and the random-time ensemble.}
With $E$ tuned to the target eigenvalue $E_0$, the all-zero acceptance probability of
an input $\ket{\psi}=\sum_k c_k\ket{E_k}$ after $M$ cycles, averaged over the
Gaussian time sets, is
\begin{equation}
  P_M=p+\sum_{k\neq0}\abs{c_k}^2\,q_k^{M},
  \qquad p=\abs{c_0}^2 ,
  \label{eq:PM}
\end{equation}
with $q_k=\tfrac12\big(1+e^{-(E_k-E_0)^2\sigma^2/2}\big)$ the Gaussian-averaged
single-cycle factor of Eq.~\eqref{eq:rodeo2}.
The target weight $p$ is preserved at every round while off-resonant components
decay geometrically, so $P_M\to p$.  Two fidelities must be distinguished for a
randomized filter.  For one fixed time set $\{t_k\}$ the post-selected state has
a fidelity $\mathcal F(\{t_k\})$ and an acceptance probability $P(\{t_k\})$; the
fidelity of the \emph{ensemble of accepted runs}, which is what repeated
executions on a device deliver, is the acceptance-weighted average
$\mathcal F_M=\langle P\,\mathcal F\rangle/\langle P\rangle$, and for it the
spectral decomposition gives exactly
\begin{equation}
  1-\mathcal{F}_M=\frac{\sum_{k\neq0}\abs{c_k}^2 q_k^{M}}
                       {p+\sum_{k\neq0}\abs{c_k}^2 q_k^{M}}
  \le\frac{(1-p)\,q_{\max}^{M}}{p+(1-p)\,q_{\max}^{M}},
  \label{eq:fid}
\end{equation}
with $q_{\max}=\tfrac12\big(1+e^{-\Delta_{\min}^2\sigma^2/2}\big)$ and
$\Delta_{\min}$ the smallest excitation energy carrying weight.  The
familiar estimate $1-\mathcal F_M\simeq\tfrac{1-p}{p}2^{-M}$ is the limit
$\Delta_{\min}\sigma\gg1$ of Eq.~\eqref{eq:fid}; for the pairing spectra used
below ($\Delta_{\min}\sigma\ge5$) the two are indistinguishable, and we plot
Eq.~\eqref{eq:fid} as the analytic expectation of the ensemble average.  The
unweighted mean of $\mathcal F(\{t_k\})$ over time sets, which some of our
figures report, differs from $\mathcal F_M$ by less than one percent for inputs
with $p>0.99$ and by a few percent for $p\simeq0.87$; both are quoted where the
difference is visible.  Reaching an ensemble fidelity $1-\delta$ needs
$M\gtrsim\log_{1/q_{\max}}[(1-p)/(p\,\delta)]$ cycles, and since the expected
number of repetitions per accepted run is $1/P_M\to1/p$, the total
controlled-evolution effort scales as $O\!\big(\abs{\log\delta}/(p\,\epsilon)\big)$
at energy resolution $\epsilon$~\cite{Choi2021}.  The filter is efficient
whenever $p$ is close to one, whatever supplies that overlap --- resolution
refinement, a variational circuit or a cheap correlated reference --- which is
why the comparison of Sec.~\ref{sec:depth} is between \emph{inputs} to the
same filter.

% ===========================================================================
\section{Results}
% ===========================================================================

\subsection{Coarse solver: UCCSD versus exact references}
\label{sec:res-uccsd}
Table~\ref{tab:uccsd} lists the ground-state energies at $g=1$ for $N=4$ as the basis
grows from $k=2$ to $k=8$.  The gate-level UCCSD circuit reproduces the operator-level
(statevector) UCCSD to machine precision and stays variational ($E_{\rm UCCSD}\ge
E_{\rm FCI}$), tracking FCI to $\sim10^{-3}$--$10^{-2}$, whereas CCD over-binds.
Figure~\ref{fig:methods} compares all solvers (exact, CCD, statevector and gate-level
UCCSD, and a hardware-efficient ansatz~\cite{Kandala2017}) as a function of $k$.  The hardware-efficient
ansatz is exact at $k=2$ and reaches FCI at $k=3$ only after a multi-start
optimization, and struggles at $k=4$ --- the symmetry-blind landscape that motivates
the physically structured UCCSD circuit.  Figure~\ref{fig:scaling} shows that the same
UCCSD solver remains a tight variational bound all the way to $k=8$ (sixteen qubits),
at the cost of a CNOT count that grows steeply ($1312$ CNOTs already at $k=4$); the
compiled circuits and their depths are analyzed in Sec.~\ref{sec:circuits}.
Two remarks keep this table in proportion.  First, an energy error is not an
infidelity: for a normalized state $\expval{H}-E_0=\sum_{j>0}\abs{c_j}^2(E_j-E_0)
\ge\Delta\,(1-\mathcal F)$, so with the physical gap $\Delta=1.823$ the $k=4$
energy error $1.44\times10^{-3}$ already caps the infidelity at
$7.9\times10^{-4}$; the actual value, obtained from the optimized state, is
$1-\mathcal F_{\rm UCCSD}=1.78\times10^{-4}$, a factor eight below the energy
error, and it is this infidelity that every comparison below uses.  Second, the
scan in $k$ at fixed $N=4$ is a basis-size study --- the fixed-$N$ sector grows
as $k^{4}$ and the seniority-zero sector as $k^{2}$ --- and it concerns the
UCCSD solver alone, not the complete pipeline; it is not evidence of favorable
many-particle scaling at fixed filling.

\begin{table}[t]
\caption{\label{tab:uccsd}Ground-state energies at $g=1$, $N=4$, for the
constant-pairing Hamiltonian.  FCI is the exact $N$-sector result; CCD is
coupled-cluster doubles; UCCSD is the gate-level circuit of Sec.~\ref{sec:uccsd}.
``\#CNOT'' is the two-qubit gate count of the UCCSD state preparation.}
\begin{ruledtabular}
\begin{tabular}{cccccc}
$k$ & qubits & FCI & CCD & UCCSD & \#CNOT\\
\colrule
2 & 4  & $1.000000$ & $1.000000$ & $1.000000$ & 0\\
3 & 6  & $0.794697$ & $0.794697$ & $0.794697$ & 272\\
4 & 8  & $0.635548$ & $0.630443$ & $0.636987$ & 1312\\
6 & 12 & $0.372664$ & $0.350058$ & $0.378719$ & --\\
8 & 16 & $0.145559$ & $0.097263$ & $0.157856$ & --\\
\end{tabular}
\end{ruledtabular}
\end{table}

\begin{figure*}[t]
\centering
\includegraphics[width=\linewidth]{figs/fig4.png}
\caption{Ground-state energy at $N=4$, $g=1$ obtained with all solvers (left) and
their deviation from FCI (right), versus the number of levels $k$.  The structured
UCCSD circuits track FCI from above; CCD over-binds; the hardware-efficient ansatz
matches FCI for $k\le3$ but loses accuracy at $k=4$.}
\label{fig:methods}
\end{figure*}

\begin{figure*}[t]
\centering
\includegraphics[width=\linewidth]{figs/fig3.png}
\caption{Scalability of the Trotterized UCCSD--VQE: ground-state energy (left) and
deviation from FCI (right) from $k=2$ to $k=8$ ($N=4$, $g=1$).  The same solver
remains variational and tight as the basis --- and hence the qubit count --- grows.}
\label{fig:scaling}
\end{figure*}

\subsection{Resolution refinement: overlap and energy}
\label{sec:res-refine}
At $N=4$ the low-resolution ($k=2$) space is a single filled determinant, so the
prolonged state is the high-resolution HF reference, with first-guess overlaps
$0.917$ for refinement to $k=3$ and $0.868$ to $k=4$.  Adiabatic evolution along the
gapped path of Eq.~\eqref{eq:path} drives both quantities toward their exact targets:
the overlap rises to $\simeq0.999$ and the energy falls from $E_{\rm HF}=1$ to the FCI
value from above --- $0.7964$ for $k=3$ (FCI $0.794697$) and $0.6397$ for $k=4$ (FCI
$0.635548$) at $T=30$ (Fig.~\ref{fig:fourparticle}).  The instantaneous spectrum stays
open along the whole path (Fig.~\ref{fig:refine}, right), a necessary condition
for the favorable $T\sim 1/\Delta E$ timescale discussed in
Sec.~\ref{sec:refine}; the larger $k=4$ space starts lower in overlap and
converges more slowly, as expected from the system-size dependence of the
refinement.

\begin{figure*}[t]
\centering
\includegraphics[width=\linewidth]{figs/fig5.png}
\caption{Resolution refinement $k=2\to k=4$.  Left: overlap of the evolved state with
the exact high-resolution ground state versus the total adiabatic time $T$, for $N=2$
and $N=4$.  Right: the instantaneous spectrum along the interpolation parameter
$s=t/T$ stays smooth and gapped, the condition for efficient refinement.}
\label{fig:refine}
\end{figure*}

\begin{figure*}[t]
\centering
\includegraphics[width=\linewidth]{figs/fig6.png}
\caption{Four-particle refinement up to $k=4$: overlap (left) and physical energy
$\expval{\Hhigh}{\Phi(T)}$ (right) versus $T$ for $k_{\rm high}=3,4$.  The energy
approaches the FCI value (dotted) from above while the overlap approaches unity, both
on a $\sim1/\Delta E$ timescale.}
\label{fig:fourparticle}
\end{figure*}

\subsection{A nontrivial coarse stage: the $k=3\to4$ chain}
\label{sec:res-k34}
At $N=4$ the $k=2$ sector is a single determinant, so in the chain above the
coarse ``solver'' is gate-free and the pipeline's first stage does no work.  To
exercise the pipeline with a working coarse stage we start instead from the
gate-level UCCSD state at $k=3$ --- a genuine $272$-CNOT variational preparation,
indistinguishable from FCI at this resolution (Table~\ref{tab:uccsd}) --- prolong
$3\to4$, and refine along the same path, with the shift $\mu$ now referenced to
the \emph{variational} coarse energy $\expval{\Hlow}{\psi_{\rm low}}$, the only
quantity a device can know.  Figure~\ref{fig:chains} compares the two chains.
The better coarse start raises the prolonged overlap from $0.868$ to $0.964$ and
reaches the target overlap $0.99$ at $T^{*}=9$ instead of $T^{*}=16$.  Whether
this saves gates is decided by the embedded coarse term of
Eq.~\eqref{eq:Aterm}, which is not free when the coarse space is nontrivial.
For the $2\to4$ chain $A$ is a single four-qubit conditional phase ($14$ CNOTs
per application, Sec.~\ref{sec:trotterized}); for the $3\to4$ chain it is the
complete $k=3$ Hamiltonian conditioned on the two appended qubits being empty,
about $350$ CNOTs per application, more than the $296$ CNOTs of the
high-resolution pairing step itself.  At the fixed step density of
Table~\ref{tab:depth} ($80$ second-order steps per $T=30$) the $2\to4$ chain
needs $43$ steps for $T^{*}=16$, i.e.\ $2\times43\times(296+14)\simeq2.7\times10^{4}$
CNOTs, and the $3\to4$ chain $24$ steps for $T^{*}=9$, i.e.\
$2\times24\times(296+350)\simeq3.1\times10^{4}$ CNOTs plus the $272$ of the coarse
circuit.  The shorter adiabatic time of the better start is thus offset by the
cost of carrying the coarse Hamiltonian along the path, and the two chains cost
the same within $15\%$.  Starting higher in the hierarchy remains the right
choice when the prolonged overlap of the trivial start is poor (the
strong-coupling regime of Sec.~\ref{sec:res-gscan}), but it is not a gate
saving in itself.

\begin{figure*}[t]
\centering
\includegraphics[width=\linewidth]{figs/fig10.png}
\caption{Refinement chains to $k=4$ at $N=4$, $g=1$, starting from the gate-free
$k=2$ determinant (circles) and from the $272$-CNOT gate-level UCCSD state at
$k=3$ (squares).  Left: overlap with the exact ground state versus adiabatic
time $T$; right: the same on a logarithmic scale.  The $3\to4$ chain starts at
overlap $0.964$ (versus $0.868$) and reaches $0.99$ in roughly half the
adiabatic time.}
\label{fig:chains}
\end{figure*}

\subsection{Coupling scan and the path gap}
\label{sec:res-gscan}
All results above sit at $g=1$; Fig.~\ref{fig:gapscan} probes the refinement
across couplings and tests the $T\sim1/\Delta E$ claim directly.  Scanning
$g=0.3$--$4$ at fixed $\mu=E_{\rm low}+0.6$ (left panel), the minimum
instantaneous gap along the path stays open and nearly constant --- it is set by
the shift itself, $\Delta E_{\min}\simeq\mu-E_{\rm low}$, dropping only to
$0.57$ at $g=4$ --- while the time to reach overlap $0.99$ grows from
$T^{*}=1$ at $g=0.3$ to $T^{*}=65$ at $g=4$.  The growth tracks the initial
overlap of the prolonged state, which falls from $0.99$ to $0.38$ over the same
range: at fixed gap, a poorer starting state simply has more rearrangement to
push through the same bottleneck.  This identifies the practical limiter of the
method at strong coupling --- not a closing gap but a degrading prolongation ---
and it is precisely the regime where a coarse start higher in the hierarchy
(Sec.~\ref{sec:res-k34}) pays off most.

Because the path gap is controlled by the shift, varying $\mu$ at fixed $g=1$
changes the minimum gap over a factor five (right panel): the extracted $T^{*}$
is well described by a line in the inverse minimum gap,
$T^{*}\approx17/\Delta E_{\min}$, which is the behavior expected of a
first-order-in-$1/\Delta E$ refinement path and incompatible with a
$1/\Delta E_{\min}^{2}$ law over this range.  The scan does not, however,
isolate the gap dependence in a strict sense: changing $\mu$ changes the
interpolation operator $H(s)$ itself, not only its gap, so the result is an
empirical trend for this family of paths in a finite system, not a test of the
general assumptions under which Ref.~\cite{Bogner2026} derives the inverse-gap
scaling.  Two gaps should also be kept apart: the shift-controlled gap at
$s=0$, $\mu-E_{\rm low}$, which is an artificial parameter, and the physical
gap of $\Hhigh$ at $s=1$ ($1.823$ at $g=1$), which the shift does not touch.
The same scan doubles as the sensitivity study for the shift: the choice
$\mu_{\rm buf}=0.6$ used throughout is conservative, and enlarging the buffer
to $1.5$ \emph{shortens} the refinement to $T^{*}=6$.  For a one-dimensional
coarse space there is no upper limit to the buffer other than the growth of
the energy scale (and hence of the Trotter error) of the path; for coarse
spaces with several states the useful initial gap saturates at the coarse
excitation gap, as noted in Sec.~\ref{sec:refine}.

\begin{figure*}[t]
\centering
\includegraphics[width=\linewidth]{figs/fig11.png}
\caption{Left: adiabatic time $T^{*}$ to reach overlap $0.99$ for the
$k=2\to4$ refinement versus pairing strength $g$ (point labels: initial overlap
of the prolonged state).  The minimum path gap stays $\simeq0.6$ throughout;
$T^{*}$ grows because the starting overlap degrades.  Right: $T^{*}$ versus the
inverse minimum path gap, varied through the shift buffer
$\mu_{\rm buf}=0.3$--$1.5$ at $g=1$; the linear fit confirms
$T^{*}\sim1/\Delta E_{\min}$.}
\label{fig:gapscan}
\end{figure*}

\subsection{Rodeo filtering of the refined state}
\label{sec:res-rodeo}
We verified the rodeo primitive in two ways: the explicit ancilla circuit agrees with
the compact filter $\tfrac12(\mathbb{1}+e^{iEt}e^{-iHt})$ to $10^{-16}$, and the
all-zero probabilities reproduce Eqs.~\eqref{eq:rodeo1}--\eqref{eq:rodeo2}, with the
Gaussian-averaged single-cycle factor matching $\tfrac12(1+e^{-(\Delta E\,\sigma)^2/2})$
to three digits.

The state fed to the filter is the resolution-refined $k=4$ state of
Sec.~\ref{sec:res-refine} at $T=25$, with overlap $p=0.9976$ with the exact ground
state; for contrast we also use the unrefined prolonged guess $\Phi_0$ ($p=0.868$).
Scanning the target energy $E$ (Fig.~\ref{fig:scan}; $M=4$ cycles, $12$ time sets
per energy) produces a peak at the ground-state energy; a coarse scan
($\sigma=2$) resolves it from the excited band, and a fine scan ($\sigma=8$)
locates it at $E=0.6352$ against the exact $E_0=0.635548$.  Starting from the
refined state, essentially all spectral weight sits in the single ground-state
peak; starting from $\Phi_0$, the small residual weight on excited states is
visible as secondary structure above the $2^{-M}$ background.  The scan is part
of the cost of the method: the fine scan alone evaluates $161$ energies with
$12$ time sets each, i.e.\ about $2\times10^{3}$ four-cycle runs, each repeated
over enough shots to estimate an acceptance probability.

Fixing $E=E_0$ and applying $M$ cycles with $\sigma=3$
(Fig.~\ref{fig:convergence}; $40$ Gaussian time sets, unweighted mean of the
post-selected fidelity), the fidelity rises geometrically and the cumulative
acceptance settles at $\approx p$, in quantitative agreement with
Eqs.~\eqref{eq:PM}--\eqref{eq:fid}.  The refined input reaches an infidelity
$1-\mathcal{F}=1.4\times10^{-3}$ after one cycle, $6.5\times10^{-4}$ after two,
$2.1\times10^{-5}$ after six and $2\times10^{-6}$ after ten; the prolonged input
reaches $1.3\times10^{-4}$ after ten cycles while its acceptance holds at
$0.868$.  The dotted lines are the analytic ensemble expectation of
Eq.~\eqref{eq:fid}; the sampled points scatter around them by up to $\pm30\%$
at the smallest infidelities, the statistical precision of $40$ time sets
(the acceptance-weighted ensemble fidelity agrees with Eq.~\eqref{eq:fid} to
better than $1\%$ for the refined input).  Each cycle buys roughly a factor two
in infidelity, so reaching $10^{-6}$ from $p=0.998$ takes ten cycles, not one
or two.  Two practical checks complete the picture.  First, the target energy
on a device is the scan estimate, not the exact eigenvalue: repeating the
filter at $E=0.6352$ instead of $E_0$ changes no digit of the infidelities
quoted above.  Second, a deliberately detuned target attenuates the ground
state as well: for $E-E_0=0.05$, $0.1$ and $0.2$ the six-cycle infidelity is
unchanged within $10\%$, $35\%$ and a factor $2.7$, while the acceptance falls
from $0.998$ to $0.971$, $0.895$ and $0.656$ --- the acceptance, not the
fidelity, is the sensitive diagnostic of a wrong target, and it then no longer
tends to the input overlap.

\begin{figure*}[t]
\centering
\includegraphics[width=\linewidth]{figs/fig7.png}
\caption{Rodeo energy scan ($M=4$).  Left: coarse scan ($\sigma=2$); the dotted
vertical lines mark exact eigenvalues and the dashed line the $2^{-M}$ background.
Right: fine scan ($\sigma=8$) on the ground-state peak, located at $E=0.6352$ (exact
$E_0=0.6355$).}
\label{fig:scan}
\end{figure*}

\begin{figure*}[t]
\centering
\includegraphics[width=\linewidth]{figs/fig8.png}
\caption{Rodeo eigenstate preparation at $E=E_0$ ($k=4$, $N=4$, $g=1$,
$\sigma=3$), averaged over $40$ Gaussian time sets, for the resolution-refined
input ($T=25$, exact slices, $p=0.9976$) and the prolonged determinant
($p=0.868$), with exact controlled evolution.  Left: unweighted mean of the
post-selected infidelity versus the number of cycles $M$ (dotted: the analytic
ensemble expectation of Eq.~\eqref{eq:fid}, which for these spectra coincides
with $\tfrac{1-p}{p}2^{-M}$).  Right: the post-selection acceptance converges to
the input overlap $p$.}
\label{fig:convergence}
\end{figure*}

% ===========================================================================
\section{Quantum circuits and depth comparison}
\label{sec:circuits}
% ===========================================================================
All quantum operations in this work are compiled to a common gate set of
single-qubit rotations and CNOTs, so the resource costs of the different routes
to the high-resolution ground state --- a full-space UCCSD--VQE, the
refinement-plus-filter pipeline, and the filter applied directly to cheaper
inputs --- can be compared gate by gate.  Every number quoted below is produced
in the accompanying notebook.  Two levels of simulation must be distinguished.
The Pauli-rotation primitive, the UCCSD excitation circuits, the
hardware-efficient ansatz and the ancilla-based rodeo cycle are simulated
\emph{gate by gate} on the full $2k$-qubit register and verified against exact
exponentials.  The product-formula evolutions of Sec.~\ref{sec:trotterized},
by contrast, are simulated \emph{block by block}: each term of the Hamiltonian
splitting is a set of mutually commuting Pauli strings whose exponential is
realized exactly by the corresponding rotation--CNOT block, and we apply that
exact block exponential within the conserved $N$-particle sector.  This
reproduces faithfully the state that the compiled circuit would produce, and
the CNOT counts are those of the compiled blocks, but it does not execute
every individual rotation and CNOT; we therefore quote CNOT \emph{counts}
rather than scheduled circuit depths, which depend on the connectivity and
parallelism of a device.

\subsection{Circuit building blocks}
\label{sec:blocks}
The elementary primitive is the Pauli rotation $e^{-i\frac{\phi}{2}P}$ for a
Pauli string $P=\bigotimes_q\sigma_q$~\cite{NielsenChuang}: each active qubit is rotated into the $Z$
eigenbasis ($H$ for $X$, $HS^\dagger$ for $Y$), a CNOT ladder accumulates the
parity onto the last active qubit, a single $R_z(\phi)$ is applied, and the
ladder and basis changes are uncomputed.  A string of weight $w$ therefore costs
$2(w-1)$ CNOTs and one parametrized rotation; the circuit reproduces the exact
exponential to $3\times10^{-16}$.

\paragraph{UCCSD preparation circuit.}
Through the Jordan--Wigner map a single excitation compiles into $2$ Pauli
rotations whose weight grows with the distance between the spin orbitals
($8$ CNOTs for adjacent levels, $16$ for the $0\!\to\!4$ single at $k=3$), and a
generic double excitation into $8$ rotations ($48$--$64$ CNOTs).  The doubles
that dominate the pairing problem are special: for a time-reversed pair moving
between levels $p$ and $q$ the operators act on the adjacent qubits
$(2p,2p{+}1)$ and $(2q,2q{+}1)$, so the internal Jordan--Wigner $Z$ strings
cancel and all $8$ Pauli strings have weight $4$ --- $48$ CNOTs per pair-hopping
term \emph{independent of the level separation}.  The full state preparation at
$N=4$, $g=1$ then costs $8$ excitation operators, $40$ rotations and $272$ CNOTs
at $k=3$, and $26$ operators, $160$ rotations and $1312$ CNOTs at $k=4$
(Table~\ref{tab:uccsd}); at $k=2$ the ansatz is empty (the sector is a single
determinant) and the preparation is gate-free.

Figure~\ref{fig:circ-vqe} shows both levels of this hierarchy.  Panel (a) is
the VQE at the ansatz level: the product of excitation unitaries acts on the HF
register, the energy is estimated from repeated measurements, and a classical
optimizer updates the $P$ amplitudes $\boldsymbol\theta$ before the circuit is
run again --- it is this feedback loop, wrapped around the full-depth state
preparation, that multiplies the $1312$-CNOT cost into the total gate budget of
the traditional route.  Panel (b) is the compiled interior of a single
weight-$4$ Pauli rotation, the sub-block that occurs eight times in every
pair-hopping double and $160$ times in the $k=4$ preparation; its only
parametrized element is the central $R_z(\phi)$, with $\phi$ fixed by the
amplitude of the parent excitation operator.

\begin{figure*}[t]
\centering
\textbf{(a)}\\[0.4em]
\begin{quantikz}[ampersand replacement=\&, row sep=0.7em, column sep=1.0em]
\lstick{$\ket{\mathrm{HF}}$} \& \qwbundle{2k} \&
\gate{e^{\theta_1(\hat\tau_1-\hat\tau_1^{\dagger})}} \& \push{\;\cdots\;} \&
\gate{e^{\theta_P(\hat\tau_P-\hat\tau_P^{\dagger})}} \& \meter{\langle \hat H\rangle} \\
\end{quantikz}\\[1.4em]
\textbf{(b)}\\[0.4em]
\begin{quantikz}[ampersand replacement=\&, row sep=0.5em, column sep=0.4em]
\lstick{$2p$}     \& \gate{H} \& \ctrl{1} \&          \&          \&                  \&          \&          \& \ctrl{1} \& \gate{H} \& \\
\lstick{$2p{+}1$} \& \gate{H} \& \targ{}  \& \ctrl{1} \&          \&                  \&          \& \ctrl{1} \& \targ{}  \& \gate{H} \& \\
\lstick{$2q$}     \& \gate{H} \&          \& \targ{}  \& \ctrl{1} \&                  \& \ctrl{1} \& \targ{}  \&          \& \gate{H} \& \\
\lstick{$2q{+}1$} \& \gate{HS^{\dagger}} \&  \&       \& \targ{}  \& \gate{R_z(\phi)} \& \targ{}  \&          \&          \& \gate{SH} \& \\
\end{quantikz}
\caption{The variational quantum eigensolver as a circuit.  (a) Ansatz level:
the Trotterized UCCSD product acts on the HF register ($2k$ qubits, slashed
wire) and the energy is estimated by measuring $\langle\hat H\rangle$ (in
practice a sum of Pauli expectation values, each averaged over many shots); a
classical optimizer updates $\boldsymbol\theta$ and the entire circuit is
re-executed, once per energy or gradient evaluation.  (b) Compiled circuit of
one weight-4 Pauli rotation, here $e^{-i\frac{\phi}{2}X\otimes X\otimes X\otimes Y}$
acting on the four qubits of a pair-hopping double excitation
$(2p,2p{+}1)\!\to\!(2q,2q{+}1)$: basis changes into the $Z$ eigenbasis ($H$ for
$X$, $HS^\dagger$ for $Y$), a CNOT parity ladder, the single parametrized
$R_z(\phi)$, and the uncompute.  Eight such blocks ($48$ CNOTs) build one
double excitation; $160$ of them ($1312$ CNOTs) build the $k=4$ preparation of
Table~\ref{tab:uccsd}.}
\label{fig:circ-vqe}
\end{figure*}

\paragraph{Hardware-efficient ansatz.}
The HEA~\cite{Kandala2017} used for contrast consists of $L$ layers of
$R_y$ rotations on all $2k$ qubits interleaved with a CNOT entangler
(nearest-neighbor pairs plus a wrap-around bond), i.e.\ $2k(L{+}1)$ variational
angles and only $O(2kL)$ CNOTs: $12$ angles at $k=2$ ($L=2$) and $30$ at $k=3$
($L=4$).  Its depth is far smaller than UCCSD, but the circuit carries no
particle-number or seniority structure; a number-penalty term and multi-start
optimization are needed at $k=3$, and at $k=4$ the optimizer fails to reach the
ground state (Fig.~\ref{fig:methods}).  Depth is not the only currency:
trainability of the symmetry-blind landscape becomes the bottleneck first.

\paragraph{Rodeo cycle.}
One rodeo cycle (Fig.~\ref{fig:circ-rodeo}) adds a Hadamard on the ancilla, the
controlled evolution $e^{-iHt_k}$, the ancilla phase $P(Et_k)$ --- the element
that effectively shifts $H\to H-E$ --- a second Hadamard, and a mid-circuit
measurement; the single ancilla is reset and reused across the $M$ cycles.  The
contrast with Fig.~\ref{fig:circ-vqe}(a) is the structural point of this
section: the rodeo circuit contains \emph{no} variational parameters and closes
no optimization loop.  The times $t_k$ are drawn once from a Gaussian, the
target energy $E$ is a classical setting, and the algorithm is driven entirely
by the measurement record --- a run is accepted when all $M$ ancilla outcomes
read $0$, which applies the filter of Sec.~\ref{sec:rodeo} cycle by cycle.
Where the VQE re-executes its full-depth circuit to steer $\boldsymbol\theta$,
the rodeo circuit is executed as printed, and its only system-size-dependent
element is the controlled evolution.  The controlled evolution is
compiled from the same Pauli-rotation primitive: in each rotation sub-circuit
only the central $R_z$ needs to be controlled (if the rotation is absent the
basis changes and CNOT ladder cancel), and a controlled $R_z$ costs $2$ extra
CNOTs.  For the pairing Hamiltonian~\eqref{eq:H} a first-order Trotter step
contains the $k(k{-}1)/2$ pair-hopping terms ($8$ weight-$4$ rotations, $48$
CNOTs each), $k$ diagonal $ZZ$ rotations ($2$ CNOTs each) and $2k$ single-qubit
phases, i.e.
\begin{equation}
  N_{\rm CNOT}^{\rm step}(k)=24\,k(k-1)+2k
  \label{eq:trottercost}
\end{equation}
CNOTs uncontrolled --- $296$ at $k=4$ --- and $\approx 4.2\times10^{2}$ CNOTs per
step once the controls are added.  A rodeo cycle with $n_T$ Trotter steps
therefore costs $\approx 4.2\times10^{2}\,n_T$ CNOTs at $k=4$.

\begin{figure}[t]
\centering
\begin{quantikz}[ampersand replacement=\&, row sep=0.9em, column sep=0.6em]
\lstick{$\ket{0}$} \& \gate{H}\gategroup[2,steps=5,style={dashed,rounded
corners,inner xsep=2pt},background,label style={label
position=below,anchor=north,yshift=-0.4cm}]{repeat $M$ times, resampling $t_k$}
 \& \ctrl{1} \& \gate{P(E\,t_k)} \& \gate{H} \& \meter{} \& \\
\lstick{$\ket{\psi_{k-1}}$} \& \qwbundle{2k} \& \gate{e^{-i\hat H t_k}} \& \& \& \& \rstick{$\ket{\psi_k}$}
\end{quantikz}
\caption{One cycle of the rodeo algorithm: Hadamard on the ancilla, controlled
time evolution $e^{-i\hat Ht_k}$ of the $2k$-qubit system register with a
Gaussian-random time $t_k$, ancilla phase $P(Et_k)$, second Hadamard, and a
mid-circuit measurement of the ancilla, which is then reset and reused for the
next cycle.  A run is accepted when all $M$ outcomes read $0$; each accepted
cycle applies the filter $\frac12(\mathbb{1}+e^{iEt_k}e^{-i\hat Ht_k})$ to the
system.  Unlike the VQE of Fig.~\ref{fig:circ-vqe}(a), the circuit carries no
variational parameters and no classical optimization loop; the controlled
evolution, compiled from the same Pauli-rotation blocks as
Fig.~\ref{fig:circ-vqe}(b) with only the central $R_z$ controlled, is the sole
system-size-dependent element ($\approx4.2\times10^{2}$ CNOTs per Trotter step
at $k=4$).}
\label{fig:circ-rodeo}
\end{figure}

\subsection{The Trotterized pipeline}
\label{sec:trotterized}
The results of Sec.~\ref{sec:res-rodeo} used exact sector exponentials; on
hardware both the refinement slices and the controlled evolution are product
formulas, and their error must be quantified at the circuit level.  We split the
Hamiltonian exactly as the gate compilation does: a diagonal term $D$ (kinetic
energy plus the $p{=}q$ pairing term; single-qubit phases and $ZZ$ rotations)
and one pair-hopping term $h_{pq}$ per level pair.  The eight Pauli strings of
each $h_{pq}$ commute, so every factor $e^{-ih_{pq}\delta t}$ is realized
\emph{exactly} by its $48$-CNOT block, and the only Trotter error is the
non-commutativity between terms --- exactly as on a device.  At $k=4$, $t=1$ the
second-order (symmetric) formula reaches spectral-norm error $3.4\times10^{-2}$
with $4$ steps and $8.3\times10^{-3}$ with $8$.

\paragraph{The embedded coarse term.}
The refinement path, Eq.~\eqref{eq:path}, contains a third kind of factor,
$e^{-i\cos^2(\pi s/2)A\,\delta t}$ with $A=P(\Hlow-\mu)P^\dagger$ of
Eq.~\eqref{eq:Aterm}.  Gate-free prolongation does not make this evolution
gate-free.  For the $2\to4$ chain at $N=4$ the coarse space is the single
determinant $\ket{\rm HF}$, so $A=(E_{\rm low}-\mu)\ket{\rm HF}\!\bra{\rm HF}
=-0.6\,\ket{\rm HF}\!\bra{\rm HF}$, and within the $N=4$ sector $\ket{\rm HF}$ is
the unique state with the four appended qubits $4$--$7$ all in $\ket0$.  Its
exponential is therefore a phase conditioned on four qubits being empty,
$X^{\otimes4}\,C^{3}P(\phi)\,X^{\otimes4}$ with $\phi=0.6\cos^2(\pi s/2)\delta t$.
Expanding the projector, $\ket{1111}\!\bra{1111}=2^{-4}\sum_{S}(-1)^{|S|}Z_S$
over the $15$ nonempty subsets $S$ of the four qubits, gives a phase polynomial
of $15$ commuting $Z$-strings; ordered along a Gray code it compiles into
$14$ CNOTs and $15$ $R_z$ rotations (the naive ladder-per-string construction
costs $34$ CNOTs).  We verified that this operator reproduces the sector
matrix exponential of $A$ exactly.  For a chain that starts from a nontrivial
coarse space the term is $(\Hlow-\mu)$ on the coarse register conditioned on
the appended qubits being empty; for the $3\to4$ chain each of the $33$ Pauli
rotations of a $k=3$ step acquires a double control ($6$ CNOTs per rotation
with a three-qubit Gray-code phase polynomial), for about $350$ CNOTs per
application, which is the figure used in Sec.~\ref{sec:res-k34}.  All
refinement budgets below include the coarse term; for the $2\to4$ chain it
adds $14$ CNOTs to the $296$ of a first-order step, i.e.\ $2\times80\times14
=2\,240$ CNOTs or $5\%$ over the $T=30$ path.

Figure~\ref{fig:trotter} shows the two pipeline stages under this compilation.
For the refinement (left), $n_s=80$ second-order steps along the $T=30$ path
reproduce the exact-slice result to within its own adiabatic error, delivering
$p=0.9977$ against the exact-slice $0.9988$; first-order steps need several
hundred.  For the rodeo stage (right) the honest scheme fixes the maximum
Trotter step \emph{size} $\delta t$, so a cycle with Gaussian time $t_k$ uses
$n_k=\lceil \abs{t_k}/\delta t\rceil$ steps of size $t_k/n_k$
($\sigma\sqrt{2/\pi}/\delta t\approx13$ steps on average at $\sigma=4$,
$\delta t=0.25$).  The filter is strikingly diagnostic of its own Trotter
error: at $\delta t=1$ the corrupted evolution is simply \emph{rejected} ---
the acceptance collapses to $0.10$ and the fidelity never converges --- while at
$\delta t=0.25$ the infidelity tracks the exact-evolution curve down to
$1.2\times10^{-4}$ at $M=6$ with acceptance $0.989$.  The intermediate
$\delta t=0.5$ converges to a floor of $2\times10^{-3}$.  The picture of
Ref.~\cite{Qian2024} --- the filter projects onto the ground state of an
\emph{effective} Trotterized Hamiltonian, whose $O(\delta t^{2})$ rotation away
from the true ground state sets the attainable fidelity --- describes this
behavior, with one qualification: because the step size $t_k/n_k$ differs from
cycle to cycle, the cycles do not share a single effective Hamiltonian, and the
observed floor is the result of a mixture of nearby effective ground states
rather than a projection on one of them.  Coherent product-formula errors can
thus both shift the located energy and bias the prepared state; the rejection
of grossly inaccurate evolutions seen at $\delta t=1$ is a useful diagnostic,
not a guarantee that smaller coherent errors are removed.

\begin{figure*}[t]
\centering
\includegraphics[width=\linewidth]{figs/fig9.png}
\caption{The pipeline at the circuit level ($k=4$, $N=4$, $g=1$).  Left:
infidelity of the Trotterized refinement versus the number of product-formula
steps $n_s$ along the $T=30$ path, for first- and second-order formulas; the
dotted line is the exact-slice reference.  Right: rodeo infidelity versus the
number of cycles $M$ for a Trotterized controlled evolution with fixed step
size $\delta t$ (input: the $n_s=80$ circuit-level refined state,
$p=0.9977$).  Too-coarse steps are rejected by the filter itself (the
$\delta t=1$ acceptance collapses to $0.10$); $\delta t=0.25$ tracks the exact
evolution.}
\label{fig:trotter}
\end{figure*}

\subsection{Gate budget at matched fidelity: which input should the filter get?}
\label{sec:depth}
The question that decides the value of the refinement stage is not whether the
pipeline beats an \emph{unfiltered} variational state --- a few rodeo cycles
improve any input with large overlap --- but whether refinement followed by the
filter beats the same filter applied to the best input that is available
cheaply.  We therefore feed one and the same Trotterized filter (second order,
$\delta t=0.25$, $\sigma=4$, $\abs{t_k}$ drawn from a Gaussian, $40$ common time
sets) four inputs: (i) the prolonged $k=2$ determinant, $p=0.868$, gate-free;
(ii) the full-space gate-level UCCSD state, $p=0.99982$, $1\,312$ CNOTs;
(iii) the Trotterized refined state at $T^{*}=16$ ($43$ second-order steps,
$p=0.9922$, $2\times43\times310=26\,660$ CNOTs including the coarse term); and
(iv) the refined state at $T=30$ of Fig.~\ref{fig:trotter} ($80$ steps,
$p=0.9977$, $49\,600$ CNOTs).  For each we track the ensemble infidelity of the
accepted runs, the acceptance $P_M$, and the CNOTs per accepted preparation,
\begin{equation}
  C_M=\frac{C_{\rm prep}+\sum_{m=1}^{M}\langle P_{m-1}\rangle\,c_m}{P_M},
  \label{eq:cost}
\end{equation}
where $c_m=2\times416\,n_m$ is the cost of the $m$-th controlled cycle and the
factor $\langle P_{m-1}\rangle$ accounts for early termination of a run at the
first ancilla outcome $1$ (without early termination the numerator is simply
$C_{\rm prep}+\sum_m c_m$; the difference is below $5\%$ for all inputs).
Equation~\eqref{eq:cost} counts repeated state preparation only.  One-time
costs are reported separately: the UCCSD training took eight L-BFGS iterations,
each requiring the energy and its $26$-component gradient, i.e.\ of order $30$
distinct $1\,312$-CNOT circuits per iteration times the shots needed for the
expectation values; the energy search of the rodeo stage costs the scans of
Sec.~\ref{sec:res-rodeo}, of order $10^{3}$ four-cycle runs times their shots,
and is common to all inputs.

Figure~\ref{fig:baselines} and Table~\ref{tab:depth} give the result.  The
first observation is that every input converges to the same floor,
$1-\mathcal F\simeq1.1$--$1.5\times10^{-4}$ after six cycles, set by the
product-formula evolution at $\delta t=0.25$ and not by the input.  The second
is the ranking at fixed cycle number.  The UCCSD state, whose actual infidelity
$1.78\times10^{-4}$ is already close to that floor, reaches $1.3\times10^{-4}$
after two cycles for $2.3\times10^{4}$ CNOTs per accepted preparation; the
refined $T=30$ state reaches $5.3\times10^{-4}$ after two cycles for
$7.2\times10^{4}$ CNOTs and $1.5\times10^{-4}$ after six for $1.2\times10^{5}$;
the refined $T^{*}=16$ state is cheaper at every cycle number than the $T=30$
state at equal or better accuracy beyond $M=4$, which shows that the
refinement time of Table~\ref{tab:depth} was not chosen economically; and the
unrefined prolonged determinant, filtered directly, reaches $1.8\times10^{-3}$
after six cycles for $7.0\times10^{4}$ CNOTs, i.e.\ roughly the accuracy of the
$T^{*}=16$ chain after two cycles ($1.7\times10^{-3}$) at roughly $1.4$ times its
cost.  At matched fidelity, therefore, refinement followed by the filter is
tied within a factor $1.5$ with filtering the unrefined input, and costs three
times more than filtering the UCCSD input while delivering a lower fidelity.
The claim that the hybrid undercuts the variational preparation by orders of
magnitude, which compared the UCCSD \emph{energy error} with the pipeline
\emph{infidelity}, does not survive the comparison of like with like.

What remains in favor of the refinement stage is qualitative but real.  It
contains no variational parameters: nothing is trained at full resolution, no
expectation values are measured during state preparation, and the only
classical inputs are the shift $\mu$ and the schedule $T$, which the coupling
and shift scans of Sec.~\ref{sec:res-gscan} set from coarse-space information.
Its gate count is dominated by the refinement, which scales with the Trotter
step of the fine Hamiltonian and with $T^{*}$, and is independent of how hard
the variational landscape is.  Whether this trade --- trainability for depth
--- pays off is a question about regimes in which the full-space variational
input is poor or expensive to train.  For the pairing model with $N=4$ that
regime is not reached at any size we can validate exactly
(Sec.~\ref{sec:sizes}), and we state this as the main negative finding of the
present assessment.

\begin{figure*}[t]
\centering
\includegraphics[width=\linewidth]{figs/fig17.png}
\caption{Matched-observable comparison of inputs to the same Trotterized
rodeo filter ($k=4$, $N=4$, $g=1$; second order, $\delta t=0.25$, $\sigma=4$,
$40$ common Gaussian time sets): the prolonged $k=2$ determinant, the full-space
gate-level UCCSD state, and the Trotterized resolution-refined state at
$T^{*}=16$ and $T=30$ (the refinement cost includes the coarse term of
Eq.~\eqref{eq:Aterm}).  Left: acceptance-weighted ensemble infidelity of the
accepted runs versus CNOTs per accepted preparation, Eq.~\eqref{eq:cost}, as
the number of cycles grows from $M=0$ (unfiltered input; omitted for the
gate-free determinant) to $M=6$.  Right: the same infidelities versus $M$.  All
inputs converge to the floor set by the product formula; the UCCSD input
reaches it first and cheapest.}
\label{fig:baselines}
\end{figure*}

\begin{table*}[t]
\caption{\label{tab:depth}Gate budget per accepted preparation at $k=4$
($N=4$, $g=1$, eight system qubits) for four inputs to the same Trotterized
rodeo filter ($\delta t=0.25$, $\sigma=4$, $40$ time sets).  $p$ is the input
overlap, $C_{\rm prep}$ the CNOT count of the input (refinement: $2\,n_s$
second-order half-steps of $296+14$ CNOTs, the $14$ being the conditional
phase of the coarse term), $1-\mathcal F_M$ the acceptance-weighted infidelity
of the accepted runs, $P_M$ the acceptance and $C_M$ the CNOTs per accepted
preparation of Eq.~\eqref{eq:cost}.  One-time training (UCCSD) and the energy
scan are not included (see text).  The UCCSD energy error of $1.4\times10^{-3}$
(Table~\ref{tab:uccsd}) corresponds to the infidelity $1.8\times10^{-4}$ listed
here.}
\footnotesize
\begin{ruledtabular}
\begin{tabular}{lcc ccc ccc}
 & & & \multicolumn{3}{c}{$M=2$} & \multicolumn{3}{c}{$M=6$}\\
input to the filter & $p$ & $C_{\rm prep}$ & $1-\mathcal F_2$ & $P_2$ & $C_2$ & $1-\mathcal F_6$ & $P_6$ & $C_6$\\
\colrule
prolonged determinant ($k{=}2\!\to\!4$) & $0.868$ & $0$ & $2.9\times10^{-2}$ & $0.887$ & $2.3\times10^{4}$ & $1.8\times10^{-3}$ & $0.854$ & $7.0\times10^{4}$\\
UCCSD, full space ($26$ parameters) & $0.99982$ & $1\,312$ & $1.3\times10^{-4}$ & $0.996$ & $2.3\times10^{4}$ & $1.1\times10^{-4}$ & $0.986$ & $6.8\times10^{4}$\\
refined, $T^{*}{=}16$ ($n_s{=}43$) & $0.9922$ & $26\,660$ & $1.7\times10^{-3}$ & $0.989$ & $4.9\times10^{4}$ & $1.7\times10^{-4}$ & $0.979$ & $9.4\times10^{4}$\\
refined, $T{=}30$ ($n_s{=}80$) & $0.9977$ & $49\,600$ & $5.3\times10^{-4}$ & $0.994$ & $7.2\times10^{4}$ & $1.5\times10^{-4}$ & $0.985$ & $1.2\times10^{5}$\\
\end{tabular}
\end{ruledtabular}
\end{table*}

\subsection{Remarks on adaptive and improved VQE}
\label{sec:adapt}
The steep UCCSD gate count is not irreducible, and much recent work aims
precisely at shallower variational circuits (see Ref.~\cite{Tilly2022} for a
review).  ADAPT-VQE~\cite{Grimsley2019} grows the ansatz one operator at a
time, at each step appending the pool operator with the largest energy gradient
and re-optimizing all amplitudes; it typically reaches a target accuracy with a
small fraction of the full UCCSD operator count and correspondingly shallower
circuits.  The pairing model makes the mechanism transparent: because the
Hamiltonian conserves seniority and the reference is seniority-zero, the
gradients of all singles and of all pair-breaking doubles vanish identically,
so an adaptive selection would retain only the $\frac{N}{2}(k-\frac{N}{2})$
pair-hopping doubles --- four operators at $k=4$, $N=4$, i.e.\
$4\times48=192$ CNOTs in place of the $1312$ of the full $26$-operator ansatz,
a sevenfold compression, with the discarded amplitudes expected to vanish at
the optimum just as the singles do.  Qubit-ADAPT-VQE~\cite{Tang2021} pushes
further by drawing the pool from individual Pauli strings, removing the
eightfold string overhead per fermionic double.  These developments are
complementary to the present pipeline rather than competing with it: the
filter is agnostic to how its input is produced, so an ADAPT-style state
--- at the coarse or at the full resolution --- slots in directly, and in the
light of Table~\ref{tab:depth} a pair-adapted full-space ansatz of $192$ CNOTs
followed by two rodeo cycles would be the cheapest route of all at this system
size.  We have not implemented ADAPT-VQE here; its inclusion as a comparator is
the natural next step of the assessment.  Conversely, no variational
compression by itself supplies the geometric, optimizer-free convergence that
the rodeo filter provides once the overlap is large.

\subsection{Other model-space sizes, particle numbers, and observables}
\label{sec:sizes}
To see whether the ranking of Table~\ref{tab:depth} is an accident of the
smallest nontrivial case, Table~\ref{tab:sizes} repeats the comparison, with
exact-slice refinement and exact controlled evolution ($\sigma=3$, $40$ time
sets), for $N=4$ and $N=6$ particles on $k=4$ and $k=6$ levels (eight and
twelve qubits), always starting from the filled determinant of the smallest
coarse space ($k_{\rm low}=N/2$).  The pattern is the same throughout: the
prolonged overlap falls to $0.78$--$0.80$ at $k=6$, refinement for $T=30$
restores it to $0.997$--$0.999$, and the UCCSD input, whose infidelity grows
from $1.8\times10^{-4}$ to $1.0\times10^{-3}$ over the table, still enters the
filter with a higher overlap than the refined state and leaves it with a lower
infidelity at every cycle number.  (At $N=6$, $k=4$ the ansatz is exact: three
pairs on four levels are the particle--hole conjugate of one pair, for which
UCCSD is complete.)  The price of the UCCSD input is its circuit ---
$6\,976$ CNOTs and $92$ parameters at $k=6$, $N=4$; $9\,024$ CNOTs and $117$
parameters at $k=6$, $N=6$ --- and the training of those parameters, which the
refinement avoids at the cost of $2n_s(N_{\rm CNOT}^{\rm step}+N_{\rm CNOT}^{A})$
CNOTs with $N_{\rm CNOT}^{\rm step}=732$ at $k=6$.  This is a controlled
small-system demonstration, not a scaling study; a regime in which refinement
wins would require a Hamiltonian for which the variational input is poor, and
we do not find one within the pairing model at these sizes.

\begin{table*}[t]
\caption{\label{tab:sizes}The pipeline and the UCCSD comparator at other
sizes ($g=1$): coarse space the filled determinant at $k_{\rm low}=N/2$,
exact-slice refinement for $T=30$, exact controlled evolution with $\sigma=3$,
$40$ time sets.  ${\rm ov}_0$ and $p$ are the overlaps of the prolonged and
refined states with the exact ground state, the UCCSD columns give the
infidelity and energy error of the sector-level UCCSD state together with its
parameter and CNOT counts, and the last four columns the acceptance-weighted
infidelity after $M$ rodeo cycles from the refined and from the UCCSD input.
The register has $2k$ qubits ($8$ and $12$) and the uncontrolled first-order
Trotter step of Eq.~\eqref{eq:trottercost} costs $296$ CNOTs at $k=4$ and
$732$ at $k=6$.}
\footnotesize
\begin{ruledtabular}
\begin{tabular}{ccc cc cccc cc cc}
$(k,N)$ & dim & gap & ${\rm ov}_0$ & $p\,(T{=}30)$ & $1-\mathcal F_{\rm UCCSD}$ & $\Delta E_{\rm UCCSD}$ & params & CNOTs & \multicolumn{2}{c}{refined$+$rodeo} & \multicolumn{2}{c}{UCCSD$+$rodeo}\\
 & & & & & & & & & $M{=}2$ & $M{=}6$ & $M{=}2$ & $M{=}6$\\
\colrule
$(4,4)$ & 70  & $1.823$ & $0.868$ & $0.9988$ & $1.8\times10^{-4}$ & $1.4\times10^{-3}$ & 26  & $1\,312$ & $3.4\times10^{-4}$ & $2.3\times10^{-5}$ & $4.1\times10^{-5}$ & $1.4\times10^{-6}$\\
$(6,4)$ & 495 & $2.017$ & $0.798$ & $0.9976$ & $7.6\times10^{-4}$ & $6.1\times10^{-3}$ & 92  & $6\,976$ & $6.3\times10^{-4}$ & $4.1\times10^{-5}$ & $1.7\times10^{-4}$ & $7.6\times10^{-6}$\\
$(4,6)$ & 28  & $1.779$ & $0.894$ & $0.9991$ & $<10^{-15}$ & $<10^{-15}$ & 15  & $688$    & $2.5\times10^{-4}$ & $1.7\times10^{-5}$ & $0$ & $0$\\
$(6,6)$ & 924 & $2.052$ & $0.780$ & $0.9972$ & $1.0\times10^{-3}$ & $8.3\times10^{-3}$ & 117 & $9\,024$ & $7.2\times10^{-4}$ & $4.8\times10^{-5}$ & $2.4\times10^{-4}$ & $1.0\times10^{-5}$\\
\end{tabular}
\end{ruledtabular}
\end{table*}

\paragraph{Observables.}
A state-preparation method for nuclear structure is judged by the observables
of the state it delivers, not only by its overlap.  Table~\ref{tab:obs} lists
the level occupation numbers $n_p=\sum_\sigma\expval{n_{p\sigma}}$ and the
correlation energy $E-E_{\rm ref}$ of the states of Table~\ref{tab:depth} at
$k=4$, $N=4$, $g=1$ (for the filtered states, the acceptance-weighted average
over the $40$ time sets).  The exact ground state depletes the two occupied
levels by $0.07$ and $0.20$ particles.  The UCCSD state reproduces the
occupations to $3\times10^{-3}$ and the correlation energy to $1.4\times10^{-3}$;
the refined state at $T=30$, whose infidelity is $2.3\times10^{-3}$, deviates by
up to $0.02$ in the occupations and by $10^{-2}$ in the correlation energy,
and two rodeo cycles reduce both deviations by a factor $2$--$3$, six cycles by
a factor $4$--$5$.  Occupation numbers converge more slowly than the energy,
as expected for a quantity that is linear rather than quadratic in the error
of the state.

\begin{table}[t]
\caption{\label{tab:obs}Level occupation numbers $n_p$ and energy of the
prepared states at $k=4$, $N=4$, $g=1$ ($E_{\rm ref}=1$); the filtered states
use the Trotterized filter of Table~\ref{tab:depth} ($\delta t=0.25$) and are
acceptance-weighted averages over $40$ time sets; $a\,10^{-b}$ stands for
$a\times10^{-b}$.}
\footnotesize
\begin{ruledtabular}
\begin{tabular}{lcccccc}
state & $n_0$ & $n_1$ & $n_2$ & $n_3$ & $E-E_{\rm ref}$ & $1-\mathcal F$\\
\colrule
exact & $1.929$ & $1.803$ & $0.197$ & $0.071$ & $-0.3645$ & $0$\\
reference det. & $2$ & $2$ & $0$ & $0$ & $0$ & $1.3\,10^{-1}$\\
UCCSD & $1.930$ & $1.806$ & $0.194$ & $0.070$ & $-0.3630$ & $1.8\,10^{-4}$\\
refined ($T{=}30$) & $1.908$ & $1.795$ & $0.205$ & $0.092$ & $-0.3549$ & $2.3\,10^{-3}$\\
refined$+$rodeo $M{=}2$ & $1.920$ & $1.798$ & $0.202$ & $0.080$ & $-0.3621$ & $5.3\,10^{-4}$\\
refined$+$rodeo $M{=}6$ & $1.923$ & $1.801$ & $0.199$ & $0.077$ & $-0.3637$ & $1.5\,10^{-4}$\\
UCCSD$+$rodeo $M{=}2$ & $1.924$ & $1.801$ & $0.199$ & $0.076$ & $-0.3637$ & $1.3\,10^{-4}$\\
\end{tabular}
\end{ruledtabular}
\end{table}

% ===========================================================================
\section{Beyond seniority: adding a particle-hole interaction}
\label{sec:ph}
% ===========================================================================
Everything above rests on a Hamiltonian that is, for the purpose of testing
methods, almost too well behaved.  The pairing force of Eq.~\eqref{eq:H} moves
complete pairs and does nothing else; the seniority is conserved, the $N=4$
ground state lives in a $\binom{k}{2}$-dimensional seniority-zero space, and an
entire class of excitations is absent --- there is no matrix element between
the reference determinant and any one-particle--one-hole state.  Realistic
nuclear and electronic Hamiltonians do have that channel.  The pure pairing
model is by no means outside the reach of mean-field-based methods --- the
Richardson model has long served as a testing ground for particle-particle
RPA and its self-consistent and renormalized
variants~\cite{Hirsch2002,Dang2006} and for the quasiparticle
RPA~\cite{RingSchuck} --- but the particle-hole channel on which Tamm--Dancoff
and particle-hole RPA are built is simply closed in it.  To test the pipeline on a
Hamiltonian with the central ingredients of a realistic two-body interaction
we therefore add, following Chapter~4 of Ref.~\cite{MHJbook}, one extra term
to Eq.~\eqref{eq:H}, keeping the level scheme, the particle number $N=4$ and the
resolutions $k=2,3,4$ unchanged:
\begin{equation}
  V_{\rm ph}=-\frac{f}{2}\sum_{p,q,r=0}^{k-1}\Big(
    a^\dagger_{p\uparrow}a^\dagger_{p\downarrow}a_{q\downarrow}a_{r\uparrow}
    +{\rm h.c.}\Big).
  \label{eq:Vph}
\end{equation}
The triple sum contains three distinct processes.  The $q=r$ pieces are the
pairing term again and merely renormalize $g\to g+2f$; the $q=p\neq r$ pieces
carry a \emph{single} particle from one level to another, a $1p$--$1h$
excitation of seniority two; and the pieces with $p,q,r$ all different remove
two unpaired particles and deposit a pair elsewhere, a seniority-two $2p$--$2h$
excitation.  $V_{\rm ph}$ still conserves $N_\uparrow$ and $N_\downarrow$
separately, so $S_z$ remains a good quantum number, but the seniority does not:
at $k=4$ the $S_z=0$ sector has $36$ determinants instead of the six pair
configurations, and the compact one-qubit-per-pair encoding of the pairing
model~\cite{Lacroix2020,AychetClaisse2026} is no longer available.  We keep
the $2k$-qubit Jordan--Wigner register used throughout this work; it is the
convenient general choice, not the only one --- the conserved particle number
and $S_z$ still permit compressed encodings of the $36$-dimensional block ---
but it is the register on which the gate counts below are meaningful.  Acting
once on the reference determinant at $k=4$, the Hamiltonian now reaches eight
$1p$--$1h$ and eight seniority-two $2p$--$2h$ determinants with matrix elements
$\pm f/2$, in addition to the four seniority-zero $2p$--$2h$ states whose
coupling grows from $g/2$ to $(g+2f)/2$.  Two consequences matter for the
pipeline.  First, the filled determinant is no longer a stationary
Hartree--Fock solution: its one-particle--one-hole matrix elements
$\mel{\Phi_i^a}{H}{\Phi_0}=\pm f/2$ do not vanish (Brillouin's theorem holds
for the self-consistent determinant, which we have not re-optimized), so the
UCCSD \emph{singles}, identically zero for the pairing force, switch on and
absorb part of what an optimized reference would describe as mean-field
relaxation rather than correlation.  Second, the reference energy becomes
$E_{\rm ref}=2-g-2f$, and the quantity we call correlation energy below is
measured from this unrelaxed reference.  We use the strengths $f=0.05$, $0.2$ and $0.5$ of
Ref.~\cite{MHJbook} at $g=1$, so that every exact number below can be
cross-checked against its tables, and we apply the pipeline exactly as before
--- the only change is the Hamiltonian passed to each stage.

\subsection{Exact and variational references}
\label{sec:ph-bench}
Table~\ref{tab:ph} collects the reference energies, the exact ground-state
energies and the UCCSD energies at $k=3$ and $k=4$.  The $k=4$ exact values
reproduce Ref.~\cite{MHJbook} to all eight digits, now obtained in the full
$70$-dimensional $N=4$ sector of the eight-qubit register rather than in the
$36$-dimensional $S_z=0$ block (the ground state lies in that block for every
$f$).  As $f$ grows the correlation energy grows much faster than the
reference energy falls, from $-0.36$ at $f=0$ to $-1.69$ at $f=0.5$
(Fig.~\ref{fig:phscan}, left), the signature of a second, competing mean field.
At $k=3$ the UCCSD ansatz remains exact for every $f$, now with non-vanishing
singles amplitudes (up to $0.09$ at $f=0.5$).  At $k=4$ the ansatz energy
error grows by a factor of twenty across the scan, from $1.4\times10^{-3}$ at
$f=0$ to $3.2\times10^{-2}$ at $f=0.5$, and the infidelity of the optimized
state by the same factor, from $1.8\times10^{-4}$ to $3.7\times10^{-3}$,
tracking the weight of the exact ground state outside the seniority-zero
space, which reaches $3.9\%$ at $f=0.5$ (Fig.~\ref{fig:phscan}, right).  This
is the first qualitative difference from the pairing model: the accuracy of
the fixed single-reference ansatz, which did not depend on $g$ before, now
degrades with the strength of the pair-breaking channel.  We note that these
are the errors of the optimized Trotterized UCCSD circuit found by L-BFGS from
a zero start; they document the accuracy of this preparation, not a proven
expressivity limit of the ansatz.

\begin{table}[t]
\caption{\label{tab:ph}The pairing model with the particle-hole term of
Eq.~\eqref{eq:Vph}, $N=4$, $g=1$.  $E_{\rm ref}=1-2f$ is the energy of the
reference determinant, FCI the exact $N$-sector ground state, UCCSD the
statevector unitary coupled-cluster energy (identical to the gate-level circuit
at $k=3$), $\Delta E_{\rm UCCSD}=E_{\rm UCCSD}-E_{\rm FCI}$, $1-\mathcal F_{\rm UCCSD}$ the infidelity of the optimized UCCSD
state with respect to the exact ground state, and $w_0$ the weight of the
exact ground state in the seniority-zero subspace.  The $k=4$ FCI column agrees
with the corresponding table of Ref.~\cite{MHJbook}.}
\footnotesize
\begin{ruledtabular}
\begin{tabular}{ccccccc}
$k$ & $f$ & FCI & UCCSD & $\Delta E_{\rm UCCSD}$ & $1-\mathcal F_{\rm UCCSD}$ & $w_0$\\
\colrule
3 & 0 & $\phantom{-}0.794697$ & $\phantom{-}0.794697$ & $<10^{-14}$ & $<10^{-14}$ & $1$\\
3 & 0.05 & $\phantom{-}0.650378$ & $\phantom{-}0.650378$ & $<10^{-14}$ & $<10^{-14}$ & $0.9994$\\
3 & 0.2 & $\phantom{-}0.185803$ & $\phantom{-}0.185803$ & $<10^{-14}$ & $<10^{-14}$ & $0.9938$\\
3 & 0.5 & $-0.839538$ & $-0.839538$ & $<10^{-14}$ & $<10^{-14}$ & $0.9808$\\
\colrule
4 & 0 & $\phantom{-}0.635548$ & $\phantom{-}0.636987$ & $1.4\times10^{-3}$ & $1.8\times10^{-4}$ & $1$\\
4 & 0.05 & $\phantom{-}0.450582$ & $\phantom{-}0.453024$ & $2.4\times10^{-3}$ & $3.0\times10^{-4}$ & $0.9988$\\
4 & 0.2 & $-0.183485$ & $-0.175303$ & $8.2\times10^{-3}$ & $1.0\times10^{-3}$ & $0.9869$\\
4 & 0.5 & $-1.691737$ & $-1.659594$ & $3.2\times10^{-2}$ & $3.7\times10^{-3}$ & $0.9605$\\
\end{tabular}
\end{ruledtabular}
\end{table}

\begin{figure*}[t]
\centering
\includegraphics[width=\linewidth]{figs/fig12.png}
\caption{The pairing model with the particle-hole term, $N=4$, $g=1$, as a
function of the pair-breaking strength $f$.  Left: energy of the reference
determinant (dashed), exact ground-state energy (solid) and UCCSD energy
(dotted) at $k=3$ and $k=4$; the vertical lines mark the strengths
$f=0.05,0.2,0.5$ of Table~\ref{tab:ph}.  Right: weight of the exact ground
state outside the seniority-zero subspace (solid) and the UCCSD energy error at
$k=4$ (dotted; at $k=3$ the ansatz is exact).}
\label{fig:phscan}
\end{figure*}

\subsection{Refinement with pairs broken}
\label{sec:ph-refine}
We prolong the $k=2$ state --- at $N=4$ still a single determinant, now with
energy $E_{\rm low}=E_{\rm ref}$ --- into the $k=4$ space and evolve along the
path of Eq.~\eqref{eq:path} with the same shift rule $\mu=E_{\rm low}+0.6$, for
each $f$ (Fig.~\ref{fig:phrefine} and Table~\ref{tab:phrefine}).  Two things
change with $f$.  The target state acquires seniority-two components that the
prolonged determinant does not contain, and the effective pairing strength
$g+2f$ grows; both lower the initial overlap, from $0.868$ at $f=0$ to $0.600$
at $f=0.5$.  The minimum gap along the path, however, is still set by the shift,
$\Delta E_{\min}=0.600$ for every $f$ (the physical gap of $\Hhigh$ in fact
\emph{opens}, from $1.82$ to $2.86$), so the adiabatic time grows only because
the state has further to travel: $T^{*}$ rises from $16$ to $32$, and for every
$f$ the overlap exceeds $0.99$ with the energy converging to FCI from above.
The comparison with the coupling scan of Sec.~\ref{sec:res-gscan} is telling:
the pure pairing model at $g=2$ --- the value of $g+2f$ at $f=0.5$ --- had
${\rm ov}_0=0.617$ and $T^{*}=32$, almost exactly the numbers found here at
$f=0.5$.  The refinement cost of the seniority-breaking model is thus governed
by the renormalized pairing strength, and the genuinely new $1p$--$1h$ and
pair-breaking channels add only a small correction to it.  The $k=3\to4$ chain
of Sec.~\ref{sec:res-k34}, started from the gate-level UCCSD state whose singles
now carry non-zero angles, buys the degradation back: it starts at overlap
$0.851$ even at $f=0.5$ and reaches $0.99$ at $T^{*}=26$ instead of $32$.

\begin{table}[t]
\caption{\label{tab:phrefine}Refinement to $k=4$ ($N=4$, $g=1$,
$\mu=E_{\rm low}+0.6$) with the particle-hole term: gap of $\Hhigh$, minimum
instantaneous gap along the path, initial overlap ${\rm ov}_0$ of the prolonged
state and adiabatic time $T^{*}$ to reach overlap $0.99$ (on the $T$ grid of
Fig.~\ref{fig:gapscan}; the $T^{*}=9$ of Sec.~\ref{sec:res-k34} for the
$3\to4$ chain at $f=0$ used a finer grid), for the gate-free $2\to4$ chain and
for the $3\to4$ chain started from the gate-level UCCSD state.}
\begin{ruledtabular}
\begin{tabular}{ccccccc}
& & & \multicolumn{2}{c}{$2\to4$} & \multicolumn{2}{c}{$3\to4$}\\
$f$ & gap($\Hhigh$) & $\Delta E_{\min}$ & ${\rm ov}_0$ & $T^{*}$ & ${\rm ov}_0$ & $T^{*}$\\
\colrule
0    & $1.823$ & $0.600$ & $0.868$ & $16$ & $0.964$ & $8$\\
0.05 & $1.943$ & $0.600$ & $0.841$ & $20$ & $0.955$ & $12$\\
0.2  & $2.289$ & $0.600$ & $0.751$ & $26$ & $0.921$ & $16$\\
0.5  & $2.860$ & $0.600$ & $0.600$ & $32$ & $0.851$ & $26$\\
\end{tabular}
\end{ruledtabular}
\end{table}

\begin{figure*}[t]
\centering
\includegraphics[width=\linewidth]{figs/fig13.png}
\caption{Resolution refinement $k=2\to4$ at $N=4$, $g=1$, for the pair-breaking
strengths $f=0,0.05,0.2,0.5$.  Left: overlap with the exact ground state versus
the adiabatic time $T$ (legend: initial overlap of the prolonged determinant;
dotted line: the $0.99$ target).  Right: variational energy error
$\expval{\Hhigh}{\Phi(T)}-E_{\rm FCI}$ on a logarithmic scale.}
\label{fig:phrefine}
\end{figure*}

\subsection{Rodeo filtering with a denser spectrum}
\label{sec:ph-rodeo}
The rodeo algorithm is the stage most exposed to the extra term, because its
energy scan sees the whole spectrum of the eight-qubit Hamiltonian in the $N=4$
sector, which the particle-hole term makes both denser and lower: $19$
eigenvalues lie below $E=4$ at $f=0.5$ against $14$ at $f=0$.
Figure~\ref{fig:phscanrodeo} repeats the energy scan of Fig.~\ref{fig:scan} at
the strongest coupling, $f=0.5$, feeding the filter the resolution-refined
($T=25$) state with $p=0.980$ and, for contrast, the prolonged determinant with
$p=0.600$.  The coarse scan resolves the ground state from the excited band ---
helped by the larger gap --- and the fine scan locates the peak at $E=-1.6924$
against the exact $E_0=-1.691737$; with the refined input the excited-state
structure that the prolonged input shows above the $2^{-M}$ background is
essentially gone.  Fixing $E=E_0$ and tracking the fidelity with the number of
cycles for all four strengths (Fig.~\ref{fig:phconv}), the infidelity falls
geometrically along the $\tfrac{1-p}{p}2^{-M}$ reference and the acceptance
settles at the input overlap: after ten cycles
$1-\mathcal{F}=2.0\times10^{-6}$, $2.9\times10^{-6}$, $4.6\times10^{-6}$ and
$4.5\times10^{-5}$ for $f=0$, $0.05$, $0.2$ and $0.5$, with acceptances
$0.998$, $0.997$, $0.994$ and $0.980$.  The pair-breaking channels cost the
pipeline one order of magnitude in final infidelity at the strongest coupling
--- recoverable with three more cycles --- and nothing in acceptance beyond the
lower input overlap.

\begin{figure*}[t]
\centering
\includegraphics[width=\linewidth]{figs/fig14.png}
\caption{Rodeo energy scan ($M=4$) for the pairing model with the particle-hole
term at $f=0.5$.  Left: coarse scan ($\sigma=2$) from the prolonged determinant
($p=0.600$) and from the resolution-refined state ($p=0.980$); dotted vertical
lines mark the exact eigenvalues, the dashed line the $2^{-M}$ background.
Right: fine scan ($\sigma=8$) on the ground-state peak, located at $E=-1.6924$
(exact $E_0=-1.6917$).}
\label{fig:phscanrodeo}
\end{figure*}

\begin{figure*}[t]
\centering
\includegraphics[width=\linewidth]{figs/fig15.png}
\caption{Rodeo eigenstate preparation at $E=E_0$ with the particle-hole term,
for $f=0,0.05,0.2,0.5$, from the resolution-refined ($T=25$) input, averaged
over $40$ random time sets ($\sigma=3$).  Left: post-selected infidelity versus
the number of cycles $M$ (dotted: $\tfrac{1-p}{p}2^{-M}$).  Right: the
acceptance converges to the input overlap $p$ (dotted).}
\label{fig:phconv}
\end{figure*}

\subsection{The circuit-level budget with the particle-hole term}
\label{sec:ph-cost}
The pipeline itself is unchanged by $V_{\rm ph}$ --- prolongation, the adiabatic
path, the shift rule and the rodeo cycle were all written for the $2k$-qubit
register --- but the Hamiltonian that is Trotterized is not.
Table~\ref{tab:pauli} lists the Pauli-string content of $H$ on $2k$ qubits.
The particle-hole term multiplies the number of strings by a factor that
grows with the resolution ($1.8$, $3.5$ and $5.3$ at $k=2,3,4$; $61\to325$ at
$k=4$) and raises the largest Pauli weight from four to $2k$ ($8$ at $k=4$),
because a $1p$--$1h$ string drags a Jordan--Wigner parity chain across the
register.

To extend the circuit-level analysis of Sec.~\ref{sec:trotterized} we
decompose $V_{\rm ph}$ into exactly exponentiable blocks in the same spirit as
the pair-hopping terms.  Its $q=r$ pieces are pairing terms and only rescale
the rotation angles of the existing blocks ($g\to g+2f$), at no cost.  The
remaining pieces, one per $(p,q,r)$ with $q\neq r$, fall into two families:
$2k(k-1)$ one-particle--one-hole blocks
($a^\dagger_{p\uparrow}n_{p\downarrow}a_{r\uparrow}+{\rm h.c.}$ and the
spin-flipped partner), each a sum of four Pauli strings, and $k(k-1)(k-2)$
pair-breaking blocks (eight strings each, like a double excitation).  Within
each block the strings mutually commute, so every factor $e^{-i\hat T\delta t}$
is realized exactly by its rotations and the product-formula error again comes
only from the non-commutativity \emph{between} blocks: $55$ of them at $k=4$
instead of seven.  Two things follow, both quantified in
Fig.~\ref{fig:phtrotter} and Table~\ref{tab:phdepth}.  First, the Trotter
error at fixed step count grows only modestly --- at $t=1$, $k=4$ the
second-order error with sixteen steps is $2.1\times10^{-3}$, $3.8\times10^{-3}$
and $7.5\times10^{-3}$ for $f=0$, $0.2$, $0.5$ --- because the new blocks carry
the small coefficient $f/2$.  The Trotterized refinement reaches its
exact-slice value with the same $n_s=80$ second-order steps at $f=0.2$ (overlap
$0.9961$ against $0.9969$), and the rodeo cycle at $\delta t=0.25$ keeps an
acceptance of $0.96$ with an infidelity of $3.8\times10^{-4}$ after six cycles,
essentially the pairing-model operating point.  At $f=0.5$ the step
$\delta t=0.25$ begins to fail --- acceptance $0.84$, the symptom identified in
Sec.~\ref{sec:trotterized} --- and $\delta t=0.125$ restores it to $0.97$ with
an infidelity of $6\times10^{-4}$ against $5\times10^{-4}$ for exact evolution.
Second, the depth of one step grows ninefold: compiled block by block with
the ladder rule of Eq.~\eqref{eq:trottercost}, a first-order step costs $2728$
CNOTs instead of $296$ ($348$ rotations instead of $60$; $3424$ against $416$
once controlled), the $1p$--$1h$ family contributing $640$ and the
pair-breaking family $1792$ CNOTs; the coarse term is unchanged ($14$ CNOTs,
the coarse space is still one determinant).  The budget per accepted
preparation therefore scales almost exactly with the step cost: at the
pairing-model operating point ($n_s=80$, $M=2$, $\delta t=0.25$) it is
$6.5\times10^{5}$ CNOTs at $f=0.2$ against $7.2\times10^{4}$ for the pure pairing
force, for an infidelity of $1.0\times10^{-3}$, and at $f=0.5$ the finer rodeo
step brings the total to $8.3\times10^{5}$ CNOTs for $5.3\times10^{-3}$.  The
price of a more realistic interaction thus appears where it should, in the
depth of each Trotter step of the Hamiltonian evolution, and not in the
structure, the operating point or the convergence of the pipeline.  Measured
against the right comparator, however, the extra depth buys no accuracy
advantage: the actual infidelity of the UCCSD state is $1.0\times10^{-3}$ at
$f=0.2$ and $3.7\times10^{-3}$ at $f=0.5$, so the hybrid at $M=2$ is tied with
the \emph{unfiltered} variational state at $f=0.2$ and behind it at $f=0.5$,
and the same filter applied to the UCCSD input (last column of
Table~\ref{tab:phdepth}) reaches $5.1\times10^{-4}$ and $2.0\times10^{-3}$ at
$\delta t=0.25$, $2.6\times10^{-4}$ and $9.6\times10^{-4}$ at $\delta t=0.125$,
for $2.0$--$3.8\times10^{5}$ CNOTs per accepted preparation --- a factor two to
three below the hybrid in both accuracy and cost.  The seniority-breaking
term thus confirms the structural robustness of the pipeline and, at the same
time, the conclusion of Sec.~\ref{sec:depth}.

\begin{figure*}[t]
\centering
\includegraphics[width=\linewidth]{figs/fig16.png}
\caption{The pipeline at the circuit level with the particle-hole term
($k=4$, $N=4$, $g=1$).  Left: infidelity of the Trotterized refinement
$k=2\to4$ ($T=30$) versus the number of product-formula steps $n_s$ for
first- and second-order formulas at $f=0.2$ (solid) and $f=0.5$ (dashed);
horizontal lines mark the exact-slice refinement.  Right: rodeo infidelity
versus cycles at $f=0.5$ on the circuit-level refined input ($n_s=80$,
second order, $p=0.9785$) with the controlled evolution Trotterized at fixed
step sizes $\delta t$; the legend gives the mean number of steps per cycle and
the six-cycle acceptance.}
\label{fig:phtrotter}
\end{figure*}

\begin{table*}[t]
\caption{\label{tab:phdepth}Gate budget at $k=4$ ($N=4$, $g=1$) with the
particle-hole term, per accepted preparation, at the operating point of
Table~\ref{tab:depth} ($n_s=80$ second-order refinement steps at $T=30$, $M=2$
rodeo cycles, $\sigma=4$) for two rodeo step sizes $\delta t$.  Step costs are
the block-by-block CNOT counts (uncontrolled / controlled); the refinement
column includes the $14$-CNOT coarse term per half-step; the rodeo column uses
the mean number of Trotter steps over the two cycles, doubled for the symmetric
formula.  The last two columns give the infidelity of the hybrid and, for
comparison, the infidelity of the unfiltered UCCSD state ($1\,312$ CNOTs,
$26$ parameters, the same circuit for every $f$) and of the UCCSD state passed
through the same two-cycle filter, with its CNOTs per accepted preparation in
parentheses.  The first row is the pure pairing model of Table~\ref{tab:depth}.}
\footnotesize
\begin{ruledtabular}
\begin{tabular}{cccccccc cc}
$f$ & $\delta t$ & step (unctrl./ctrl.) & refinement & rodeo & per run & $P_M$ & per accepted & $1-\mathcal{F}$: hybrid / UCCSD & UCCSD$+$rodeo (CNOTs)\\
\colrule
$0$   & $0.25$  & $296$ / $416$   & $49\,600$  & $22\,000$ & $71\,600$ & $0.994$ & $7.2\times10^{4}$ & $5.3\times10^{-4}$ / $1.8\times10^{-4}$ & $1.3\times10^{-4}$ ($2.3\times10^{4}$)\\
$0.2$ & $0.25$  & $2728$ / $3424$ & $438\,720$ & $190\,603$ & $629\,323$ & $0.972$ & $6.5\times10^{5}$ & $1.0\times10^{-3}$ / $1.0\times10^{-3}$ & $5.1\times10^{-4}$ ($2.0\times10^{5}$)\\
$0.2$ & $0.125$ & $2728$ / $3424$ & $438\,720$ & $375\,499$ & $814\,219$ & $0.996$ & $8.2\times10^{5}$ & $8.1\times10^{-4}$ / $1.0\times10^{-3}$ & $2.6\times10^{-4}$ ($3.8\times10^{5}$)\\
$0.5$ & $0.25$  & $2728$ / $3424$ & $438\,720$ & $190\,603$ & $629\,323$ & $0.891$ & $7.1\times10^{5}$ & $6.1\times10^{-3}$ / $3.7\times10^{-3}$ & $2.0\times10^{-3}$ ($2.1\times10^{5}$)\\
$0.5$ & $0.125$ & $2728$ / $3424$ & $438\,720$ & $375\,499$ & $814\,219$ & $0.979$ & $8.3\times10^{5}$ & $5.3\times10^{-3}$ / $3.7\times10^{-3}$ & $9.6\times10^{-4}$ ($3.8\times10^{5}$)\\
\end{tabular}
\end{ruledtabular}
\end{table*}

\begin{table}[t]
\caption{\label{tab:pauli}Pauli-string content of the Hamiltonian on $2k$
qubits: number of strings and largest Pauli weight, for the pairing model and
for the pairing model with the particle-hole term ($f\neq0$; the counts do not
depend on the value of $f$).  The $k=2,3$ rows agree with Ref.~\cite{MHJbook}.}
\footnotesize
\begin{ruledtabular}
\begin{tabular}{cccc}
$k$ & qubits & pairing & pairing $+$ $V_{\rm ph}$\\
\colrule
2 & 4 & $15$ (weight $4$) & $27$ (weight $4$)\\
3 & 6 & $34$ (weight $4$) & $118$ (weight $6$)\\
4 & 8 & $61$ (weight $4$) & $325$ (weight $8$)\\
\end{tabular}
\end{ruledtabular}
\end{table}

% ===========================================================================
\section{Discussion and outlook}
\label{sec:discussion}
% ===========================================================================
The pairing model lets us exhibit the complete refinement-plus-filter chain
with every stage benchmarked against exact results: a gate-level UCCSD coarse
solver, a gate-free prolongation with a counted conditioned coarse term, a
gapped adiabatic refinement whose extracted time grows linearly with the
inverse of the shift-controlled path gap over the range studied, and a rodeo
filter whose ensemble infidelity follows Eq.~\eqref{eq:fid} --- all validated
a second time at the circuit level with product formulas.  The two algorithms
are complementary in the intended sense: refinement raises the overlap from
$0.87$ to $0.998$, and the filter removes the adiabatic tail geometrically.
The coupling scan adds a practical lesson: at strong pairing the limiter is
not the path gap but the quality of the prolonged starting state.  The
particle-hole extension of Sec.~\ref{sec:ph} shows that none of this depends
on the seniority symmetry that makes the pairing model convenient: with the
one-particle--one-hole channel open, the same stages run unchanged, the path
gap is still set by the shift, the refinement cost is governed by the
renormalized pairing strength $g+2f$, and the filter converges geometrically
for every strength studied; the cost of the more realistic interaction shows
up in the Pauli-string content of the Hamiltonian ($61\to325$ strings at $k=4$)
and hence in the depth of each Trotter step ($296\to2728$ CNOTs).

The central quantitative finding, however, is negative for the refinement
stage as a gate-saving device.  Compared on the same observable --- the
ground-state infidelity --- and with the same Trotterized filter, refinement
followed by the rodeo algorithm is tied with filtering the unrefined prolonged
determinant and is beaten, in accuracy and in cost per accepted preparation,
by filtering the full-space UCCSD state, for the pure pairing force as well as
with the particle-hole term and at every size we can validate exactly
(Tables~\ref{tab:depth}, \ref{tab:sizes} and \ref{tab:phdepth}).  The reason is
simple: in these systems a single-reference variational circuit with a few
tens of parameters already has an overlap above $0.998$, and an overlap of
$0.87$ costs the filter only a few extra cycles, whereas an adiabatic path of
length $T^{*}\sim10$--$30$ at a Trotter step of several hundred CNOTs costs
tens of thousands.  What the refinement offers in exchange is that it is
optimizer-free and measurement-free during preparation, and that its cost is
set by the Hamiltonian evolution rather than by the trainability of a
landscape.  Whether that trade pays off is a question about regimes we cannot
reach with exact validation --- larger spaces, interactions for which
single-reference ans\"atze fail, or hardware on which the training
measurements of a variational loop dominate --- and we regard the present
work as the controlled baseline against which such claims should be measured,
not as evidence for them.

\paragraph{Relation to symmetry-restoration and subspace approaches.}
The results invite a comparison with the pairing-model work of Lacroix and
collaborators summarized in the Introduction.  Both programmes start from a
cheap, imperfect trial state and sharpen it; the difference is what is
sharpened and how.  In the symmetry-restoration route the trial state breaks a
symmetry on purpose and the projector --- phase estimation on the symmetry
operator~\cite{Lacroix2020}, an oracle~\cite{RuizGuzman2023} or a
measurement-side projection~\cite{RuizGuzman2024} --- removes the unwanted
sectors, after which a Krylov or phase-estimation stage refines the
energy~\cite{RuizGuzman2022}.  Our coarse state conserves particle number and
spin by construction; what it lacks is \emph{resolution}, and the refinement
supplies it along a gapped path while the rodeo cycle filters in
\emph{energy} rather than in a symmetry quantum number.  The two filters are
structurally the same resource --- one ancilla, a controlled unitary and a
post-selected measurement --- and they compose: a symmetry-projected input
would enter our pipeline exactly as the UCCSD state does.  The quantum
subspace methods of Refs.~\cite{Zhang2025,AychetClaisse2026,BeaujeaultTaudiere2024}
are complementary in a different way: they trade coherent depth for
measurements and a classical generalized eigenvalue problem, and their
accuracy is bounded by the quality of the states that span the subspace,
whereas the rodeo filter converges geometrically in depth but requires the
controlled evolution.  Two observations of that work bear on ours, with the
limits of an energy filter kept in mind.  Ref.~\cite{Zhang2024} documents that
ADAPT-VQE with a symmetry-breaking pool can collapse to unphysical energies on
degenerate pairing systems.  A downstream energy filter can repair such an
input only insofar as the input retains a nonzero overlap with the target:
it cannot recover a state of zero overlap, cannot move weight between
particle-number sectors that the evolution conserves, and cannot separate
exactly degenerate states by energy alone, so a symmetry-broken ADAPT state
would have to be symmetry-projected before, not instead of, the filter.
Ref.~\cite{AychetClaisse2026} shows that the accuracy of the one-body Green's
function of a small superfluid system, computed with the pure pairing
Hamiltonian across the normal-to-superfluid transition driven by the pairing
strength, is controlled by the quality of the even-$N$ ground state; a
filtered state of the kind prepared here is a possible input to that
construction, and the particle-hole term of Sec.~\ref{sec:ph} --- which is
not part of the model of Ref.~\cite{AychetClaisse2026} --- is the ingredient
that forces one out of the pair-to-qubit encoding used there.

\paragraph{Hardware requirements.}
The Trotterized pipeline of Sec.~\ref{sec:trotterized} quantifies the dominant
coherent error of an ideal device, and the requirements it implies are
explicit: $2k+1$ qubits ($9$ at $k=4$, $13$ at $k=6$, $17$ at $k=8$); a
mid-circuit measurement with reset and reuse of one ancilla; arbitrary-angle
$R_z$ rotations, $160$ of them per UCCSD preparation and $60$ (pairing) to
$348$ (with $V_{\rm ph}$) per Trotter step; and a coherent budget per accepted
preparation that ranges from $2\times10^{4}$ two-qubit gates for the cheapest
route of Table~\ref{tab:depth} (UCCSD input, two cycles) to $10^{5}$ for the
refined input with six cycles, and to $10^{6}$ with the particle-hole term
(Table~\ref{tab:phdepth}), on top of the energy scan and the repetitions
needed to estimate acceptance probabilities.  At the two-qubit error rates
demonstrated in large-scale experiments today --- of order $10^{-2}$ on the
superconducting processor of Ref.~\cite{Kim2023} and $(0.8$--$2)\times10^{-3}$
on trapped-ion processors~\cite{Moses2023,Ransford2025} --- a $10^{4}$-gate
run suffers of order ten faults even on the best of these and a
$10^{5}$-gate run of order one hundred; none of these
budgets is executable with physical qubits without substantial error
mitigation, and qubit counts alone say nothing about a device's suitability.
Early logical-qubit demonstrations on neutral-atom arrays have executed
circuits on tens of logical qubits with error detection and, in part, error
correction and mid-circuit readout~\cite{Bluvstein2024,Reichardt2024}, and the
rodeo algorithm has been run in two-qubit form on trapped-ion and
superconducting hardware~\cite{BeeLindgren2022,Qian2024}; a logical-depth
implementation of the budgets above remains to be demonstrated.  A
proportionate near-term target is therefore a mitigated $k=3$ demonstration:
the $2\to3$ chain of Sec.~\ref{sec:res-refine} on seven qubits (six system
qubits plus the ancilla), whose $N=4$ coarse state is a single determinant, or
the $3\to4$ chain of Sec.~\ref{sec:res-k34} on nine qubits, with a short
refinement schedule and one or two rodeo cycles at coarse $\delta t$, i.e.\
$10^{3}$--$10^{4}$ two-qubit gates, where the collapse of the acceptance under
a corrupted evolution (Sec.~\ref{sec:trotterized}) doubles as a hardware
diagnostic.  All-to-all connectivity absorbs the Jordan--Wigner ladders
without SWAP overhead; on square or heavy-hex lattices the weight-4
pair-hopping rotations act on adjacent qubits, which limits but does not
eliminate routing costs, and the weight-$2k$ strings of $V_{\rm ph}$ do not.
The full $k=4$ budget is a workload for the first generation of logical-qubit
machines such as the neutral-atom system planned by the QuNorth
initiative~\cite{QuNorth}, and Tables~\ref{tab:depth} and \ref{tab:phdepth}
are offered as its specification.

\paragraph{Concluding remarks and perspectives.}
The present results establish, on a fully controlled example, how resolution
refinement followed by the rodeo algorithm behaves at the circuit level ---
every stage compiled, counted and validated against exact results --- and
what it is worth: it prepares the correlated ground state to the accuracy of
the Trotterized evolution without any variational optimization, at a gate
cost that, for the pairing model at these sizes, exceeds that of filtering a
cheap variational input.  Three extensions would sharpen the assessment.
First, comparators: an ADAPT-VQE or pair-adapted input, and a jointly
optimized choice of refinement time and cycle number along the cost--fidelity
curves of Fig.~\ref{fig:baselines}, with statistical uncertainties from
larger time-set ensembles.  Second, Hamiltonians for which single-reference
variational inputs fail --- shell-model interactions with competing
deformations, or the strong-coupling regime of Sec.~\ref{sec:res-gscan}
treated with a nontrivial coarse space --- which is where the trainability
argument for refinement would be tested rather than asserted.  Third,
hardware: a mitigated seven- or nine-qubit demonstration on present devices,
where the acceptance statistics of the filter double as a diagnostic of the
compiled evolution, followed by the full $k=4$ budget on an early
logical-qubit machine.  The pairing model, exactly solvable at every
resolution, remains the calibration standard for all three.

% ===========================================================================
\section*{Code availability}
The complete implementation --- pairing Hamiltonian with and without the
particle-hole term, gate-level UCCSD--VQE, prolongation, adiabatic refinement,
and the ancilla-based rodeo algorithm --- is provided as a tested Python package
and a single Jupyter notebook that reproduces every figure in this article.
Give GitHub link when everything is ready.
\begin{acknowledgments}
The work of C.G. and G.F.L. acknowledge funding from the
European Union’s Horizon Europe research and innovation program under the Marie Skłodowska-Curie Grant Agreement No.
101126636. M.H.-J. is partially supported by U.S. National Science Foundation
(NSF) Grants PHY-1404159 and PHY-2310020. 
The authors  acknowledge Anthropic’s Claude
for assistance with literature search and coding.

\end{acknowledgments}

% ===========================================================================
\begin{thebibliography}{99}
% ===========================================================================
\bibitem{Peruzzo2014}
A.~Peruzzo, J.~McClean, P.~Shadbolt, M.-H.~Yung, X.-Q.~Zhou, P.~J.~Love,
A.~Aspuru-Guzik, J.~L.~O'Brien,
\textit{A variational eigenvalue solver on a photonic quantum processor},
Nat.\ Commun.\ \textbf{5}, 4213 (2014).

\bibitem{McClean2018}
J.~R.~McClean, S.~Boixo, V.~N.~Smelyanskiy, R.~Babbush, H.~Neven,
\textit{Barren plateaus in quantum neural network training landscapes},
Nat.\ Commun.\ \textbf{9}, 4812 (2018).

\bibitem{Albash2018}
T.~Albash, D.~A.~Lidar,
\textit{Adiabatic quantum computation},
Rev.\ Mod.\ Phys.\ \textbf{90}, 015002 (2018).

\bibitem{Abrams1999}
D.~S.~Abrams, S.~Lloyd,
\textit{Quantum algorithm providing exponential speed increase for finding
eigenvalues and eigenvectors},
Phys.\ Rev.\ Lett.\ \textbf{83}, 5162 (1999).

\bibitem{Choi2021}
K.~Choi, D.~Lee, J.~Bonitati, Z.~Qian, J.~Watkins,
\textit{Rodeo algorithm for quantum computing},
Phys.\ Rev.\ Lett.\ \textbf{127}, 040505 (2021).

\bibitem{Qian2024}
Z.~Qian, J.~Watkins, G.~Given, J.~Bonitati, K.~Choi, D.~Lee,
\textit{Demonstration of the rodeo algorithm on a quantum computer},
Eur.\ Phys.\ J.\ A \textbf{60}, 151 (2024).

\bibitem{Bogner2026}
S.~Bogner, H.~Hergert, M.~Hjorth-Jensen, R.~LaRose, D.~Lee, M.~Patkowski,
\textit{Quantum state preparation with resolution refinement},
Phys.\ Lett.\ B \textbf{875}, 140363 (2026); arXiv:2511.14732.

\bibitem{Patkowski2025fusion}
M.~Patkowski, O.~Ayyildiz, M.~Kebri\v{c}, K.~L.~C.~Hunt, D.~Lee,
\textit{Hierarchical fusion method for scalable quantum eigenstate preparation},
Phys.\ Rev.\ A \textbf{113}, 052442 (2026); arXiv:2510.19039.

\bibitem{LeeReview2023}
D.~Lee,
\textit{Quantum techniques for eigenvalue problems},
Eur.\ Phys.\ J.\ A \textbf{59}, 275 (2023).

\bibitem{LeeUCC2019}
J.~Lee, W.~J.~Huggins, M.~Head-Gordon, K.~B.~Whaley,
\textit{Generalized unitary coupled cluster wavefunctions for quantum computation},
J.\ Chem.\ Theory Comput.\ \textbf{15}, 311 (2019).

\bibitem{Kandala2017}
A.~Kandala, A.~Mezzacapo, K.~Temme, M.~Takita, M.~Brink, J.~M.~Chow,
J.~M.~Gambetta,
\textit{Hardware-efficient variational quantum eigensolver for small molecules
and quantum magnets},
Nature \textbf{549}, 242 (2017).

\bibitem{Grimsley2019}
H.~R.~Grimsley, S.~E.~Economou, E.~Barnes, N.~J.~Mayhall,
\textit{An adaptive variational algorithm for exact molecular simulations on a
quantum computer},
Nat.\ Commun.\ \textbf{10}, 3007 (2019).

\bibitem{Tang2021}
H.~L.~Tang, V.~O.~Shkolnikov, G.~S.~Barron, H.~R.~Grimsley, N.~J.~Mayhall,
E.~Barnes, S.~E.~Economou,
\textit{Qubit-ADAPT-VQE: an adaptive algorithm for constructing
hardware-efficient ans\"atze on a quantum processor},
PRX Quantum \textbf{2}, 020310 (2021).

\bibitem{Tilly2022}
J.~Tilly, H.~Chen, S.~Cao, D.~Picozzi, K.~Setia, Y.~Li, E.~Grant, L.~Wossnig,
I.~Rungger, G.~H.~Booth, J.~Tennyson,
\textit{The variational quantum eigensolver: a review of methods and best
practices},
Phys.\ Rep.\ \textbf{986}, 1 (2022).

\bibitem{Trotter1959}
H.~F.~Trotter,
\textit{On the product of semi-groups of operators},
Proc.\ Am.\ Math.\ Soc.\ \textbf{10}, 545 (1959).

\bibitem{Suzuki1976}
M.~Suzuki,
\textit{Generalized Trotter's formula and systematic approximants of exponential
operators},
Commun.\ Math.\ Phys.\ \textbf{51}, 183 (1976).

\bibitem{Childs2021}
A.~M.~Childs, Y.~Su, M.~C.~Tran, N.~Wiebe, S.~Zhu,
\textit{Theory of Trotter error with commutator scaling},
Phys.\ Rev.\ X \textbf{11}, 011020 (2021).

\bibitem{LNP936}
M.~Hjorth-Jensen, M.~P.~Lombardo, U.~van~Kolck (eds.),
\textit{An Advanced Course in Computational Nuclear Physics},
Lecture Notes in Physics \textbf{936}, Springer (2017).

\bibitem{MHJbook}
M.~Hjorth-Jensen,
\textit{Quantum Mechanics for Many-Particle Systems},
lecture notes, University of Oslo (2026), Chapter~4.

\bibitem{Lacroix2020}
D.~Lacroix,
\textit{Symmetry assisted preparation of entangled many-body states on a quantum computer},
Phys.\ Rev.\ Lett.\ \textbf{125}, 230502 (2020).

\bibitem{RuizGuzman2021}
E.~A.~Ruiz~Guzman, D.~Lacroix,
\textit{Calculation of generating function in many-body systems with quantum computers: technical challenges and use in hybrid quantum-classical methods},
arXiv:2104.08181 (2021).

\bibitem{Siwach2021}
P.~Siwach, D.~Lacroix,
\textit{Filtering states with total spin on a quantum computer},
Phys.\ Rev.\ A \textbf{104}, 062435 (2021).

\bibitem{RuizGuzman2022}
E.~A.~Ruiz~Guzman, D.~Lacroix,
\textit{Accessing ground-state and excited-state energies in a many-body system after symmetry restoration using quantum computers},
Phys.\ Rev.\ C \textbf{105}, 024324 (2022).

\bibitem{Lacroix2023EPJA}
D.~Lacroix, E.~A.~Ruiz~Guzman, P.~Siwach,
\textit{Symmetry breaking/symmetry preserving circuits and symmetry restoration on quantum computers: A quantum many-body perspective},
Eur.\ Phys.\ J.\ A \textbf{59}, 3 (2023).

\bibitem{RuizGuzman2023}
E.~A.~Ruiz~Guzman, D.~Lacroix,
\textit{Restoring broken symmetries using quantum search ``oracles''},
Phys.\ Rev.\ C \textbf{107}, 034310 (2023).

\bibitem{Ayral2023}
T.~Ayral, P.~Besserve, D.~Lacroix, E.~A.~Ruiz~Guzman,
\textit{Quantum computing with and for many-body physics},
Eur.\ Phys.\ J.\ A \textbf{59}, 227 (2023).

\bibitem{BeaujeaultTaudiere2024}
Y.~Beaujeault-Taudi\`ere, D.~Lacroix,
\textit{Solving the Lipkin model using quantum computers with two qubits only with a hybrid quantum-classical technique based on the generator coordinate method},
Phys.\ Rev.\ C \textbf{109}, 024327 (2024).

\bibitem{RuizGuzman2024}
E.~A.~Ruiz~Guzman, D.~Lacroix,
\textit{Restoring symmetries in quantum computing using Classical Shadows},
Eur.\ Phys.\ J.\ A \textbf{60}, 112 (2024).

\bibitem{Zhang2025}
J.~Zhang, D.~Lacroix,
\textit{Excited states from ADAPT-VQE convergence path in many-body problems: application to nuclear pairing problem and H$_4$ molecule dissociation},
Phys.\ Lett.\ B \textbf{869}, 139841 (2025).

\bibitem{AychetClaisse2026}
S.~Aychet-Claisse, D.~Lacroix, V.~Som\`a, J.~Zhang,
\textit{Quantum simulations of Green's functions for small superfluid systems},
Phys.\ Rev.\ C \textbf{113}, 044324 (2026); arXiv:2509.02272.

\bibitem{RingSchuck}
P.~Ring, P.~Schuck,
\textit{The Nuclear Many-Body Problem},
Springer, Berlin (1980).

\bibitem{Hirsch2002}
J.~G.~Hirsch, A.~Mariano, J.~Dukelsky, P.~Schuck,
\textit{Fully self-consistent RPA description of the many level pairing model},
Ann.\ Phys.\ (N.Y.) \textbf{296}, 187 (2002).

\bibitem{Dang2006}
N.~Dinh~Dang,
\textit{Consistency of particle-particle random-phase approximation and its
renormalizations},
Phys.\ Rev.\ C \textbf{74}, 034326 (2006).

\bibitem{NielsenChuang}
M.~A.~Nielsen, I.~L.~Chuang,
\textit{Quantum Computation and Quantum Information},
Cambridge University Press (2010).

\bibitem{Richardson1963}
R.~W.~Richardson,
\textit{A restricted class of exact eigenstates of the pairing-force Hamiltonian},
Phys.\ Lett.\ \textbf{3}, 277 (1963).

\bibitem{RichardsonSherman1964}
R.~W.~Richardson, N.~Sherman,
\textit{Exact eigenstates of the pairing-force Hamiltonian},
Nucl.\ Phys.\ \textbf{52}, 221 (1964).

\bibitem{Dukelsky2004}
J.~Dukelsky, S.~Pittel, G.~Sierra,
\textit{Colloquium: Exactly solvable Richardson-Gaudin models for many-body
quantum systems},
Rev.\ Mod.\ Phys.\ \textbf{76}, 643 (2004).

\bibitem{Ovrum2007}
E.~Ovrum, M.~Hjorth-Jensen,
\textit{Quantum computation algorithm for many-body studies},
arXiv:0705.1928 (2007).

\bibitem{Hlatshwayo2022}
M.~Q.~Hlatshwayo, Y.~Zhang, H.~Wibowo, R.~LaRose, D.~Lacroix, E.~Litvinova,
\textit{Simulating excited states of the Lipkin model on a quantum computer},
Phys.\ Rev.\ C \textbf{106}, 024319 (2022).

\bibitem{Stetcu2022}
I.~Stetcu, A.~Baroni, J.~Carlson,
\textit{Variational approaches to constructing the many-body nuclear ground state
for quantum computing},
Phys.\ Rev.\ C \textbf{105}, 064308 (2022).

\bibitem{Stetcu2023}
I.~Stetcu, A.~Baroni, J.~Carlson,
\textit{Projection algorithm for state preparation on quantum computers},
Phys.\ Rev.\ C \textbf{108}, L031306 (2023).

\bibitem{PerezObiol2023}
A.~P\'erez-Obiol, A.~M.~Romero, J.~Men\'endez, A.~Rios, A.~Garc\'ia-S\'aez,
B.~Juli\'a-D\'iaz,
\textit{Nuclear shell-model simulation in digital quantum computers},
Sci.\ Rep.\ \textbf{13}, 12291 (2023).

\bibitem{Jiang2023}
C.~Jiang, J.~Pei,
\textit{Quantum computing of the pairing Hamiltonian at finite temperature},
Phys.\ Rev.\ C \textbf{107}, 044308 (2023).

\bibitem{Zhang2024}
J.~Zhang, D.~Lacroix, Y.~Beaujeault-Taudi\`ere,
\textit{Neutron-proton pairing correlations described on quantum computers},
Phys.\ Rev.\ C \textbf{110}, 064320 (2024).

\bibitem{Lacroix2023}
E.~A.~Ruiz~Guzman,
\textit{Symmetry breaking and restoration for many-body problems treated on
quantum computers},
PhD thesis, Universit\'e Paris-Saclay (2023); arXiv:2310.17996.

\bibitem{Cohen2023}
T.~D.~Cohen, H.~Oh,
\textit{Optimizing the rodeo projection algorithm},
Phys.\ Rev.\ A \textbf{108}, 032422 (2023).

\bibitem{BeeLindgren2022}
M.~Bee-Lindgren, Z.~Qian, M.~DeCross, N.~C.~Brown, C.~N.~Gilbreth, J.~Watkins,
X.~Zhang, D.~Lee,
\textit{Controlled gate networks: theory and application to eigenvalue estimation},
Eur.\ Phys.\ J.\ A \textbf{61}, 263 (2025); arXiv:2208.13557.

\bibitem{Patkowski2026}
M.~Patkowski, O.~Ayyildiz, K.~Hunt, N.~Jansen, D.~Lee,
\textit{Improved rodeo algorithm performance for spectral functions and state
preparation},
arXiv:2602.05978 (2026).

\bibitem{LinTong2020}
L.~Lin, Y.~Tong,
\textit{Near-optimal ground state preparation},
Quantum \textbf{4}, 372 (2020).

\bibitem{Dong2022}
Y.~Dong, L.~Lin, Y.~Tong,
\textit{Ground-state preparation and energy estimation on early fault-tolerant
quantum computers via quantum eigenvalue transformation of unitary matrices},
PRX Quantum \textbf{3}, 040305 (2022).

\bibitem{Larocca2025}
M.~Larocca, S.~Thanasilp, S.~Wang, K.~Sharma, J.~Biamonte, P.~J.~Coles,
L.~Cincio, J.~R.~McClean, Z.~Holmes, M.~Cerezo,
\textit{Barren plateaus in variational quantum computing},
Nat.\ Rev.\ Phys.\ \textbf{7}, 174 (2025).

\bibitem{Beck2023}
D.~Beck \textit{et al.},
\textit{Quantum information science and technology for nuclear physics.
Input into U.S.\ long-range planning, 2023},
arXiv:2303.00113 (2023).

\bibitem{Bauer2023}
C.~W.~Bauer, Z.~Davoudi, N.~Klco, M.~J.~Savage,
\textit{Quantum simulation of fundamental particles and forces},
Nat.\ Rev.\ Phys.\ \textbf{5}, 420 (2023).

\bibitem{QuNorth}
QuNorth, \textit{Magne: a level-2 neutral-atom quantum computer},
\url{https://qunorth.com} (accessed July 2026).

\bibitem{Kim2023}
Y.~Kim, A.~Eddins, S.~Anand, K.~X.~Wei, E.~van~den~Berg, S.~Rosenblatt,
H.~Nayfeh, Y.~Wu, M.~Zaletel, K.~Temme, A.~Kandala,
\textit{Evidence for the utility of quantum computing before fault tolerance},
Nature \textbf{618}, 500 (2023).

\bibitem{Moses2023}
S.~A.~Moses \textit{et al.},
\textit{A race-track trapped-ion quantum processor},
Phys.\ Rev.\ X \textbf{13}, 041052 (2023).

\bibitem{Ransford2025}
A.~Ransford \textit{et al.},
\textit{Helios: a 98-qubit trapped-ion quantum computer},
arXiv:2511.05465 (2025).

\bibitem{Bluvstein2024}
D.~Bluvstein \textit{et al.},
\textit{Logical quantum processor based on reconfigurable atom arrays},
Nature \textbf{626}, 58 (2024).

\bibitem{Reichardt2024}
B.~W.~Reichardt \textit{et al.},
\textit{Logical computation demonstrated with a neutral atom quantum processor},
arXiv:2411.11822 (2024).
\end{thebibliography}

\end{document}
